\chapter{补充材料与证明}\label{app:proof}
\section{变分视角}\label{app-sec:var-proof}


\subsection{定理~\ref{thm:equiv-marginal-kl}：边缘与条件 KL 最小化的等价性}\label{subsec:proof-equiv-marginal-kl}
\paragraph{\textit{证明。}} \textbfs{\Cref{eq:kl-matching} 的推导。}

我们首先展开右端的期望：
\begin{align*}
    &\mathbb{E}_{p(\rvx_0, \rvx_i)}\left[
    \mathcal{D}_{\text{KL}} \big(p(\rvx_{i-1}|\rvx_i, \rvx_0) \| p_\phi(\rvx_{i-1}|\rvx_i)\big)
    \right]
    \\
    =& \int \int p(\rvx_0, \rvx_i)   
    \mathcal{D}_{\text{KL}} \big(p(\rvx_{i-1}|\rvx_i, \rvx_0) \| p_\phi(\rvx_{i-1}|\rvx_i)\big) 
    \diff\rvx_0 \diff\rvx_i.
\end{align*}
根据 KL 散度的定义，
\begin{align*}
    &\mathcal{D}_{\text{KL}} \big(p(\rvx_{i-1}|\rvx_i, \rvx_0) \| p_\phi(\rvx_{i-1}|\rvx_i)\big)
    = \int p(\rvx_{i-1}|\rvx_i, \rvx_0) 
    \log \frac{p(\rvx_{i-1}|\rvx_i, \rvx_0)}{p_\phi(\rvx_{i-1}|\rvx_i)} 
    \diff\rvx_{i-1}.
\end{align*}
将其代入期望，得
\begin{align*}
    &\int \int \int p(\rvx_0, \rvx_i)  p(\rvx_{i-1}|\rvx_i, \rvx_0) 
    \log \frac{p(\rvx_{i-1}|\rvx_i, \rvx_0)}{p_\phi(\rvx_{i-1}|\rvx_i)} 
    \diff\rvx_{i-1} \diff\rvx_0 \diff\rvx_i.
\end{align*}
利用概率的链式法则，
\begin{align*}
    p(\rvx_0, \rvx_i) 
    &= p(\rvx_i)   p(\rvx_0|\rvx_i),
\end{align*}
我们将积分改写为
\begin{align*}
    &\int p(\rvx_i) \int p(\rvx_0|\rvx_i) \int p(\rvx_{i-1}|\rvx_i, \rvx_0)
    \log \frac{p(\rvx_{i-1}|\rvx_i, \rvx_0)}{p_\phi(\rvx_{i-1}|\rvx_i)} 
    \diff\rvx_{i-1} \diff\rvx_0 \diff\rvx_i.
\end{align*}
这使我们可以将期望表示为嵌套形式：
\begin{align*}
    &\mathbb{E}_{p(\rvx_i)} \left[
    \mathbb{E}_{p(\rvx_0|\rvx_i)} \left[
    \mathbb{E}_{p(\rvx_{i-1}|\rvx_i, \rvx_0)} \left[
    \log \frac{p(\rvx_{i-1}|\rvx_i, \rvx_0)}{p_\phi(\rvx_{i-1}|\rvx_i)}
    \right]
    \right]
    \right].
\end{align*}

接下来，我们应用对数的分解：
\begin{align*}
    \log \frac{p(\rvx_{i-1}|\rvx_i, \rvx_0)}{p_\phi(\rvx_{i-1}|\rvx_i)}
    =\log \frac{p(\rvx_{i-1}|\rvx_i, \rvx_0)}{p(\rvx_{i-1}|\rvx_i)} 
    + \log \frac{p(\rvx_{i-1}|\rvx_i)}{p_\phi(\rvx_{i-1}|\rvx_i)}.
\end{align*}
将其代回期望得到两项：
\begin{align*}
    &\mathbb{E}_{p(\rvx_i)} \left[
    \mathbb{E}_{p(\rvx_0|\rvx_i)} \left[
    \mathbb{E}_{p(\rvx_{i-1}|\rvx_i, \rvx_0)} \left[
    \log \frac{p(\rvx_{i-1}|\rvx_i, \rvx_0)}{p(\rvx_{i-1}|\rvx_i)}
    \right]
    \right]
    \right]
    \\
    &+
    \mathbb{E}_{p(\rvx_i)} \left[
    \mathbb{E}_{p(\rvx_0|\rvx_i)} \left[
    \mathbb{E}_{p(\rvx_{i-1}|\rvx_i, \rvx_0)} \left[
    \log \frac{p(\rvx_{i-1}|\rvx_i)}{p_\phi(\rvx_{i-1}|\rvx_i)}
    \right]
    \right]
    \right].
\end{align*}
由于第二个对数项不依赖于 $\rvx_0$，根据全概率律
\begin{align*}
    \mathbb{E}_{p(\rvx_0|\rvx_i)} \left[
    \mathbb{E}_{p(\rvx_{i-1}|\rvx_i, \rvx_0)} \left[
    \log \frac{p(\rvx_{i-1}|\rvx_i)}{p_\phi(\rvx_{i-1}|\rvx_i)}
    \right]
    \right]
    = \mathbb{E}_{p(\rvx_{i-1}|\rvx_i)} \left[
    \log \frac{p(\rvx_{i-1}|\rvx_i)}{p_\phi(\rvx_{i-1}|\rvx_i)}
    \right].
\end{align*}
类似地，第一项是 KL 散度
\begin{align*}
    \mathbb{E}_{p(\rvx_0|\rvx_i)} \left[
    \mathcal{D}_{\text{KL}} \big(p(\rvx_{i-1}|\rvx_i, \rvx_0)  \|  p(\rvx_{i-1}|\rvx_i)\big)
    \right].
\end{align*}
将所有内容综合起来，我们得到如下分解：
\begin{align*}
    &\mathbb{E}_{p(\rvx_0, \rvx_i)} \left[
    \mathcal{D}_{\text{KL}} \big(p(\rvx_{i-1}|\rvx_i, \rvx_0)  \|  p_\phi(\rvx_{i-1}|\rvx_i)\big)
    \right]
    \\
    =& \mathbb{E}_{p(\rvx_i)} \left[
    \mathbb{E}_{p(\rvx_0|\rvx_i)} \left[
    \mathcal{D}_{\text{KL}} \big(p(\rvx_{i-1}|\rvx_i, \rvx_0)  \|  p(\rvx_{i-1}|\rvx_i)\big)
    \right]
    \right]
    \\
    &+
    \mathbb{E}_{p(\rvx_i)} \left[
    \mathcal{D}_{\text{KL}} \big(p(\rvx_{i-1}|\rvx_i)  \|  p_\phi(\rvx_{i-1}|\rvx_i)\big)
    \right].
\end{align*}

\textbfs{最优性的证明。} 要证：
\[
p^*(\rvx_{i-1}|\rvx_i)
=
p(\rvx_{i-1}|\rvx_i)
=
\mathbb{E}_{p(\rvx_0|\rvx_i)}
\left[
p(\rvx_{i-1}|\rvx_i,\rvx_0)
\right],
\qquad
\rvx_i \sim p_i.
\]
第一个等式由以下事实得到：在假设参数化具有足够表达能力的前提下，当 $p^* = p$ 时 KL 散度 $\mathcal{D}_{\mathrm{KL}}(p \| p_{\bm{\phi}})$ 取得最小值。第二个等式直接由全概率律得到。
\hfill$\blacksquare$



\subsection{定理~\ref{thm:ddpm-elbo}：扩散模型的 ELBO}\label{app:vlb-proof}

\paragraph{\textit{证明。}} 为简化记号，我们记
\[
\rvx_{1:L}:=(\rvx_1,\ldots,\rvx_L).
\]

\textbfs{第一步：应用 Jensen 不等式。}
边缘对数似然为
\[
\log p_{\bm{\phi}}(\rvx_0)
=
\log \int p_{\bm{\phi}}(\rvx_0,\rvx_{1:L})\,\diff \rvx_{1:L},
\]
其中联合生成分布为
\[
p_{\bm{\phi}}(\rvx_0,\rvx_{1:L})
=
p_{\text{prior}}(\rvx_L)\prod_{i=1}^L p_{\bm{\phi}}(\rvx_{i-1}| \rvx_i).
\]

我们引入前向加噪分布
\[
p(\rvx_{1:L}| \rvx_0)
=
\prod_{i=1}^L p(\rvx_i| \rvx_{i-1}),
\]
并改写为
\[
\log p_{\bm{\phi}}(\rvx_0)
=
\log \int
p(\rvx_{1:L}| \rvx_0)
\frac{p_{\bm{\phi}}(\rvx_0,\rvx_{1:L})}{p(\rvx_{1:L}| \rvx_0)}
\,\diff \rvx_{1:L}.
\]
应用 Jensen 不等式得到
\[
\log p_{\bm{\phi}}(\rvx_0)
\ge
\mathbb{E}_{p(\rvx_{1:L}| \rvx_0)}
\left[
\log
\frac{p_{\bm{\phi}}(\rvx_0,\rvx_{1:L})}{p(\rvx_{1:L}| \rvx_0)}
\right]
=: \mathcal{L}_{\text{ELBO}}(\rvx_0;\bm{\phi}),
\]
因此
\[
-\log p_{\bm{\phi}}(\rvx_0)
\le
-\mathcal{L}_{\text{ELBO}}(\rvx_0;\bm{\phi}).
\]

\textbfs{第二步：展开 ELBO。}
将各分解式代入 ELBO，得
\begin{align*}
\mathcal{L}_{\text{ELBO}}(\rvx_0;\bm{\phi})
&=
\mathbb{E}_{p(\rvx_{1:L}| \rvx_0)}
\Bigg[
\log p_{\text{prior}}(\rvx_L)
+
\sum_{i=1}^L \log p_{\bm{\phi}}(\rvx_{i-1}| \rvx_i)
-
\sum_{i=1}^L \log p(\rvx_i| \rvx_{i-1})
\Bigg].
\end{align*}
因此，
\begin{align*}
-\mathcal{L}_{\text{ELBO}}(\rvx_0;\bm{\phi})
&=
\mathbb{E}_{p(\rvx_{1:L}| \rvx_0)}
\Bigg[
\log \frac{p(\rvx_L| \rvx_0)}{p_{\text{prior}}(\rvx_L)}
+
\sum_{i=2}^L
\log \frac{p(\rvx_{i-1}| \rvx_i,\rvx_0)}{p_{\bm{\phi}}(\rvx_{i-1}| \rvx_i)}
-
\log p_{\bm{\phi}}(\rvx_0| \rvx_1)
\Bigg].
\end{align*}
这里我们用到了如下恒等式
\[
p(\rvx_{1:L}| \rvx_0)
=
p(\rvx_L| \rvx_0)\prod_{i=2}^L p(\rvx_{i-1}| \rvx_i,\rvx_0),
\]
它由前向过程的马尔可夫结构得到。

现在根据依赖关系将各项分组：
\begin{align*}
-\mathcal{L}_{\text{ELBO}}(\rvx_0;\bm{\phi})
&=
\underbrace{
\mathbb{E}_{p(\rvx_L| \rvx_0)}
\left[
\log \frac{p(\rvx_L| \rvx_0)}{p_{\text{prior}}(\rvx_L)}
\right]
}_{\mathcal{L}_{\text{prior}}(\rvx_0)}
+
\underbrace{
\mathbb{E}_{p(\rvx_1| \rvx_0)}
\left[
-\log p_{\bm{\phi}}(\rvx_0| \rvx_1)
\right]
}_{\mathcal{L}_{\text{recon.}}(\rvx_0;\bm{\phi})}
\\
&\quad+
\underbrace{
\sum_{i=2}^L
\mathbb{E}_{p(\rvx_i| \rvx_0)}
\left[
\mathcal{D}_{\mathrm{KL}}
\big(
p(\rvx_{i-1}| \rvx_i,\rvx_0)
\|
p_{\bm{\phi}}(\rvx_{i-1}| \rvx_i)
\big)
\right]
}_{\mathcal{L}_{\text{diffusion}}(\rvx_0;\bm{\phi})}.
\end{align*}
这就证明了所声称的分解。
\hfill$\blacksquare$


\clearpage
\newpage

\section{分数视角}\label{app-sec:score-proof}

\subsection{命题~\ref{sm-trace}：通过分部积分实现可解的分数匹配}\label{app-sec:sm-trace}
\paragraph{\textit{证明。}}

 \textbfs{展开 $\mathcal{L}_{\text{SM}}(\bm{\phi})$。 }
        让我们展开期望内部的平方差：
        \begin{align*}
        \mathcal{L}_{\text{SM}}(\bm{\phi}) &= \frac{1}{2}\mathbb{E}_{\rvx \sim p_{\text{data}}(\rvx)} \left[ \norm{\rvs_{\bm{\phi}}(\rvx)}_2^2 - 2 \langle \rvs_{\bm{\phi}}(\rvx), \rvs(\rvx) \rangle + \norm{\rvs(\rvx)}_2^2 \right] \\
        &= \frac{1}{2}\mathbb{E}_{\rvx \sim p_{\text{data}}(\rvx)} \left[ \norm{\rvs_{\bm{\phi}}(\rvx)}_2^2 \right] - \mathbb{E}_{\rvx \sim p_{\text{data}}(\rvx)} \left[ \langle \rvs_{\bm{\phi}}(\rvx), \rvs(\rvx) \rangle \right] \\
        &\quad + \frac{1}{2}\mathbb{E}_{\rvx \sim p_{\text{data}}(\rvx)} \left[ \norm{\rvs(\rvx)}_2^2 \right].
        \end{align*}
        我们现在关注交叉项：
        \[
        \mathbb{E}_{\rvx \sim p_{\text{data}}(\rvx)} \left[ \langle \rvs_{\bm{\phi}}(\rvx), \rvs(\rvx) \rangle \right].
        \]
        利用如下事实
        \[
        \nabla_{\rvx} \log p_{\text{data}}(\rvx) = \frac{\nabla_{\rvx}p_{\text{data}}(\rvx)}{p_{\text{data}}(\rvx)},
        \]
        并假设 $p_{\text{data}}(\rvx)$ 不为零（例如在其支撑集上），交叉项变为：
        \begin{align*}
            \mathbb{E}_{\rvx \sim p_{\text{data}}(\rvx)} \left[ \langle \rvs_{\bm{\phi}}(\rvx), \rvs(\rvx) \rangle \right] 
             &= \int \rvs_{\bm{\phi}}(\rvx)^\top  \nabla_{\rvx} \log p_{\text{data}}(\rvx) p_{\text{data}}(\rvx) \diff \rvx \\
             &= \int \rvs_{\bm{\phi}}(\rvx)^\top  \nabla_{\rvx}p_{\text{data}}(\rvx) \diff \rvx \\
             &= \sum_{i=1}^{D} \int  s_{\bm{\phi}}^{(i)}(\rvx) \partial_{x_i}p_{\text{data}}(\rvx) \diff \rvx,
        \end{align*}
        其中 $s_{\bm{\phi}}^{(i)}(\rvx)$ 是分数函数的第 $i$ 个分量
        \[
        \rvs_{\bm{\phi}} = \left( s_{\bm{\phi}}^{(1)}, s_{\bm{\phi}}^{(2)}, \dots, s_{\bm{\phi}}^{(D)} \right).
        \]
        \proofparagraph{分部积分。}
        我们使用如下分部积分公式~\citep{evans10}，它由标准微积分推导而来：
        \begin{lemma*}
            设 $u, v$ 为半径 $R>0$ 的球 $\mathbb{B}(\bm{0}, R) \subset \mathbb{R}^D$ 上的可微实值函数。则对 $i = 1, \ldots, D$，下式成立：
        \[
        \int_{\mathbb{B}(\bm{0}, R)} u   \partial_{x_i} v   \diff \rvx = -\int_{\mathbb{B}(\bm{0}, R)} v   \partial_{x_i} u   \diff \rvx + \int_{\partial \mathbb{B}(\bm{0}, R)} u v \nu_i   \diff S,
        \]
        其中 $\nu = (\nu_1, \ldots, \nu_D)$ 是边界 $\partial \mathbb{B}(\bm{0}, R)$——半径为 $R>0$ 的球面——的单位外法向量，$\diff S$ 是 $\partial \mathbb{B}(\bm{0}, R)$ 上的曲面测度。
        \end{lemma*}
        我们对所有 $i = 1, \dots, D$ 将该公式应用于 $u(\rvx) := s_{\bm{\phi}}^{(i)}(\rvx)$ 和 $v(\rvx) = p_{\text{data}}(\rvx)$，假设
        \[
        |u(\rvx) v(\rvx)| \to 0 \quad \text{as } R \to \infty.
        \]
        将结果对所有 $i = 1, \dots, D$ 求和，得：
        \begin{align*}
            \mathbb{E}_{\rvx \sim p_{\text{data}}(\rvx)} \left[ \langle \rvs_{\bm{\phi}}(\rvx), \rvs(\rvx) \rangle \right] 
             &= - \sum_{i=1}^{D} \int  \partial_{x_i}s_{\bm{\phi}}^{(i)}(\rvx) p_{\text{data}}(\rvx) \diff \rvx \\
             &= - \mathbb{E}_{\rvx \sim p_{\text{data}}(\rvx)} \left[\Tr\left(\nabla_{\rvx}\rvs_{\bm{\phi}}(\rvx)\right) \right].
        \end{align*}
        综合所有结果，我们有：
        \begin{align*}
        \mathcal{L}_{\text{SM}}(\bm{\phi}) &= \underbrace{\mathbb{E}_{\rvx \sim p_{\text{data}}(\rvx)} \left[\Tr\left(\nabla_{\rvx}\rvs_{\bm{\phi}}(\rvx)\right) + \frac{1}{2}  \norm{\rvs_{\bm{\phi}}(\rvx)}_2^2 \right]}_{\widetilde{\mathcal{L}}_{\text{SM}}(\bm{\phi})} \\
        &\quad + \underbrace{\frac{1}{2}\mathbb{E}_{\rvx \sim p_{\text{data}}(\rvx)} \left[ \norm{\rvs(\rvx)}_2^2 \right]}_{=:C},
        \end{align*}
        其中 $C$ 仅依赖于分布 $p_{\text{data}}$，证明完毕。
\hfill$\blacksquare$

\subsection{定理~\ref{thm:sm-dsm}：SM 与 DSM 最小化的等价性}\label{app-sec:sm-dsm}
\paragraph{\textit{证明。}} 展开 $\mathcal{L}_{\text{SM}}(\bm{\phi}; \sigma)$ 和 $\mathcal{L}_{\text{DSM}}(\bm{\phi}; \sigma)$，得：
\begin{align*}
 &\mathcal{L}_{\text{SM}}(\bm{\phi}; \sigma) = \frac{1}{2} \mathbb{E}_{\tilde{\rvx} \sim p_\sigma(\tilde{\rvx})} 
 \begin{aligned}[t]
         \Big[ 
    \norm{\rvs_{\bm{\phi}}(\tilde{\rvx}; \sigma)}_2^2 
    &- 2 \rvs_{\bm{\phi}}(\tilde{\rvx}; \sigma)^\top  \nabla_{\tilde{\rvx}} \log p_\sigma(\tilde{\rvx}) \\
    &+ \norm{\nabla_{\tilde{\rvx}} \log p_\sigma(\tilde{\rvx})}_2^2 
    \Big], 
 \end{aligned}
\\
&\mathcal{L}_{\text{DSM}}(\bm{\phi}; \sigma) = \frac{1}{2} \mathbb{E}_{ p_{\text{data}}(\rvx)p_{\sigma}(\tilde{\rvx} | \rvx)}
 \begin{aligned}[t]
         \Big[
    \norm{\rvs_{\bm{\phi}}(\tilde{\rvx}; \sigma)}_2^2 
    &- 2 \rvs_{\bm{\phi}}(\tilde{\rvx}; \sigma)^\top \nabla_{\tilde{\rvx}} \log p_{\sigma}(\tilde{\rvx} | \rvx) \\
    &+ \norm{\nabla_{\tilde{\rvx}} \log p_{\sigma}(\tilde{\rvx} | \rvx)}_2^2 
    \Big].
 \end{aligned}
\end{align*}
两式相减得：
\begin{align*}
 &\mathcal{L}_{\text{SM}}(\bm{\phi}; \sigma) -\mathcal{L}_{\text{DSM}}(\bm{\phi}; \sigma) 
 \\= \frac{1}{2} &\Bigg( \mathbb{E}_{\tilde{\rvx} \sim p_\sigma(\tilde{\rvx})} \norm{\rvs_{\bm{\phi}}(\tilde{\rvx}; \sigma)}_2^2 - \mathbb{E}_{ p_{\text{data}}(\rvx)p_{\sigma}(\tilde{\rvx} | \rvx)} \norm{\rvs_{\bm{\phi}}(\tilde{\rvx}; \sigma)}_2^2\Bigg)
 \\ \quad - &  
 \begin{aligned}[t]
 \Bigg(
      \mathbb{E}_{\tilde{\rvx} \sim p_\sigma(\tilde{\rvx})}  &\left[ \rvs_{\bm{\phi}}(\tilde{\rvx}; \sigma)^\top  \nabla_{\tilde{\rvx}} \log p_\sigma(\tilde{\rvx}) \right] \\ 
      &- \mathbb{E}_{ p_{\text{data}}(\rvx)p_{\sigma}(\tilde{\rvx} | \rvx)}\left[ \rvs_{\bm{\phi}}(\tilde{\rvx}; \sigma)^\top  \nabla_{\tilde{\rvx}} \log p_\sigma(\tilde{\rvx}| \rvx)\right]
      \Bigg)
 \end{aligned}
 \\ \quad +  \frac{1}{2}&\Bigg(\mathbb{E}_{\tilde{\rvx} \sim p_\sigma(\tilde{\rvx})}  \norm{\nabla_{\tilde{\rvx}} \log p_\sigma(\tilde{\rvx})}_2^2  -  \mathbb{E}_{ p_{\text{data}}(\rvx)p_{\sigma}(\tilde{\rvx} | \rvx)} \norm{\nabla_{\tilde{\rvx}} \log p_{\sigma}(\tilde{\rvx} | \rvx)}_2^2   \Bigg).
\end{align*}
接下来，我们逐一处理这三项。对于第一项，由于 $p_\sigma(\tilde{\rvx}) = \int p_{\sigma}(\tilde{\rvx} | \rvx) p_{\text{data}}(\rvx) \diff\rvx$，我们可以将其改写为：
\begin{align*}
    \mathbb{E}_{\tilde{\rvx} \sim p_\sigma(\tilde{\rvx})} \norm{\rvs_{\bm{\phi}}(\tilde{\rvx}; \sigma)}_2^2 
    &= \int \Big( \int p_{\sigma}(\tilde{\rvx} | \rvx) p_{\text{data}}(\rvx)  \diff\rvx \Big) \norm{\rvs_{\bm{\phi}}(\tilde{\rvx}; \sigma)}_2^2 \diff \tilde{\rvx} 
    \\ &= \int p_{\text{data}}(\rvx) \int p_{\sigma}(\tilde{\rvx} | \rvx) \norm{\rvs_{\bm{\phi}}(\tilde{\rvx}; \sigma)}_2^2 \diff \tilde{\rvx} \diff\rvx
    \\ &= \mathbb{E}_{ p_{\text{data}}(\rvx)p_{\sigma}(\tilde{\rvx} | \rvx)} \norm{\rvs_{\bm{\phi}}(\tilde{\rvx}; \sigma)}_2^2.
\end{align*}
因此，第一项为零。对于第二项：
\begin{align}\label{eq:dsm-sec-term}
\begin{aligned}
          &\mathbb{E}_{\tilde{\rvx} \sim p_\sigma(\tilde{\rvx})}\left[ \rvs_{\bm{\phi}}(\tilde{\rvx}; \sigma)^\top  \nabla_{\tilde{\rvx}} \log p_\sigma(\tilde{\rvx}) \right] 
      \\=&\int  p_\sigma(\tilde{\rvx}) \rvs_{\bm{\phi}}(\tilde{\rvx}; \sigma)^\top  \frac{\nabla_{\tilde{\rvx}} p_\sigma(\tilde{\rvx})}{p_\sigma(\tilde{\rvx})}  \diff\tilde\rvx
        \\=&\int   \rvs_{\bm{\phi}}(\tilde{\rvx}; \sigma)^\top  \nabla_{\tilde{\rvx}} \int p_{\sigma}(\tilde{\rvx} | \rvx) p_{\text{data}}(\rvx) \diff\rvx \diff\tilde\rvx
    \\=&\int \int  \rvs_{\bm{\phi}}(\tilde{\rvx}; \sigma)^\top    \nabla_{\tilde{\rvx}}p_{\sigma}(\tilde{\rvx} | \rvx) p_{\text{data}}(\rvx) \diff\tilde\rvx \diff\rvx 
      \\=  &\mathbb{E}_{ p_{\text{data}}(\rvx)p_{\sigma}(\tilde{\rvx} | \rvx)}\left[ \rvs_{\bm{\phi}}(\tilde{\rvx}; \sigma)^\top  \nabla_{\tilde{\rvx}} \log p_\sigma(\tilde{\rvx}| \rvx)\right].
\end{aligned}
\end{align}
因此，它也为零。对于第三项，注意：
\[
C:=\frac{1}{2}\Bigg(\mathbb{E}_{\tilde{\rvx} \sim p_\sigma(\tilde{\rvx})}  \norm{\nabla_{\tilde{\rvx}} \log p_\sigma(\tilde{\rvx})}_2^2  -  \mathbb{E}_{ p_{\text{data}}(\rvx)p_{\sigma}(\tilde{\rvx} | \rvx)} \norm{\nabla_{\tilde{\rvx}} \log p_{\sigma}(\tilde{\rvx} | \rvx)}_2^2   \Bigg)
\]
仅依赖于 $p_{\text{data}}(\rvx)$ 和 $p_{\sigma}(\tilde{\rvx}|\rvx)$，因此关于 $\bm{\phi}$ 是常数。
\hfill$\blacksquare$

\subsection{引理~\ref{tweedie}：Tweedie 公式}\label{app-sec:tweedie}
我们首先陈述 Tweedie 公式的一种更一般的形式，它考虑了依赖于时间的高斯扰动，并在下面给出其证明。


\paragraph{具有时间依赖参数的 Tweedie 恒等式。}
设 $\rvx_t \sim \mathcal{N}\big(\cdot; \alpha_t\rvx_0, \sigma_t^2\rmI\big)$ 为高斯随机向量。则 Tweedie 公式表明
\begin{align*}
\alpha_t \mathbb{E}_{\mathbf{x}_0 \sim p(\mathbf{x}_0|\rvx_t)}[\mathbf{x}_0|\rvx_t] = \rvx_t + \sigma_t^2 \nabla_{\rvx_t} \log p_t(\rvx_t),
\end{align*}
其中期望是在给定观测 $\rvx_t$ 下 $\mathbf{x}_0$ 的后验分布 $p(\mathbf{x}_0|\rvx_t)$ 上取的，$p_t(\rvx_t)$ 是 $\rvx_t$ 的边缘密度。


\paragraph{\textit{证明。}}


\paragraph{边缘密度及其分数。}
我们回忆 $\rvx_t$ 的边缘密度由下式给出
\begin{align*}
p_t(\rvx_t) = \int p_t(\rvx_t|\rvx_0)  p_0(\rvx_0)  \diff\rvx_0.
\end{align*}
我们现在计算分数函数：
\begin{align*}
\nabla_{\rvx_t} \log p_t(\rvx_t) 
= \frac{\nabla_{\rvx_t} p_t(\rvx_t)}{p_t(\rvx_t)} = \frac{1}{p_t(\rvx_t)} \int \nabla_{\rvx_t} p_t(\rvx_t|\rvx_0)  p_0(\rvx_0)  \diff\rvx_0.
\end{align*}
因此我们需要计算条件密度的梯度。

\paragraph{条件密度的梯度与整理。}
条件高斯密度的梯度为：
\begin{align*}
\nabla_{\rvx_t} p_t(\rvx_t|\rvx_0) 
= -p_t(\rvx_t|\rvx_0) \cdot \sigma_t^{-2} (\rvx_t - \alpha_t \rvx_0).
\end{align*}
将其代入前一个表达式，得：
\begin{align*}
\nabla_{\rvx_t} p_t(\rvx_t) 
&= \int \nabla_{\rvx_t} p_t(\rvx_t|\rvx_0)  p_0(\rvx_0)  \diff\rvx_0 \\
&= -\sigma_t^{-2} \int (\rvx_t - \alpha_t \rvx_0)  p_t(\rvx_t|\rvx_0)  p_0(\rvx_0)  \diff\rvx_0 \\
&= -\sigma_t^{-2} \int (\rvx_t - \alpha_t \rvx_0)  p(\rvx_0|\rvx_t)  p_t(\rvx_t)  \diff\rvx_0 \\
&= -p_t(\rvx_t)  \sigma_t^{-2} \left(\rvx_t - \alpha_t  \mathbb{E}_{p(\rvx_0|\rvx_t)}[\rvx_0\vert \rvx_t] \right).
\end{align*}
两边同时除以 $p_t(\rvx_t)$，得：
\begin{align*}
\nabla_{\rvx_t} \log p_t(\rvx_t) 
= -\sigma_t^{-2} \left(\rvx_t - \alpha_t  \mathbb{E}_{p(\rvx_0|\rvx_t)}[\rvx_0\vert \rvx_t] \right).
\end{align*}
整理得：
\begin{align*}
\rvx_t + \sigma_t^{2} \nabla_{\rvx_t} \log p_t(\rvx_t) 
= \alpha_t  \mathbb{E}_{p(\rvx_0|\rvx_t)}[\rvx_0\vert \rvx_t].
\end{align*}
推导完毕。

\subsection{Stein 恒等式与代理 SURE 目标}\label{app-sec:stein-identity}

\paragraph{Stein 恒等式。} Stein 恒等式是一种分部积分技巧，它将未知密度下的期望转化为可观测函数及其导数的期望，从而消去配分函数，使得无需计算未知密度或配分函数即可获得无偏的、可解的目标函数和检验。我们从最简单的一维情形出发，然后将其推广到证明 SURE 代理损失所需的形式。
\subparagraph{一维标准正态情形。}
若 $z \sim \mathcal{N}(0,1)$ 且 $f$ 具有适当的衰减性，则 Stein 恒等式表明：
\[
\E[f'(z)]  =  \E[z f(z)] .
\]
记 $\phi(z) := \tfrac{1}{\sqrt{2\pi}}  e^{-z^2/2}$ 为一维标准正态密度。
证明通过分部积分完成，利用 $\phi'(z) = -z\phi(z)$ 以及消失的边界项。
为精确说明，我们计算
\[
\E[f'(z)] 
= \int_{-\infty}^{\infty} f'(z)  \phi(z)  \diff z.
\]
通过分部积分，取 $u=\phi(z)$、$\diff v=f'(z)  \diff z$，得
\[
\int f'(z)  \phi(z)  \diff z
= \Big[ f(z)  \phi(z) \Big]_{-\infty}^{\infty}
- \int f(z)  \phi'(z)  \diff z.
\]
由于 $\phi'(z) = -z\phi(z)$ 且当 $|z|\to\infty$ 时 $f(z)\phi(z)\to 0$（衰减条件），边界项消失，故
\[
\E[f'(z)]  =  \int f(z)  z  \phi(z)  \diff z  =  \E[z f(z)].
\]
这就完成了推导，并证明了 Stein 恒等式的一维情形。


\subparagraph{多元标准正态情形。}
若 $\rvz \sim \mathcal{N}(\mathbf{0},\rmI_D)$ 且 $g:\R^D \to \R$，则 Stein 恒等式为
\[
\E[\nabla g(\rvz)]  =  \E[\rvz g(\rvz)].
\]
等价地，对于 $\mathbf{u}:\R^D \to \R^D$，
\begin{align}\label{eq:stein-multi}
    \E[ \nabla_{\tilde{\rvx}} \cdot\mathbf{u}(\rvz)]  =  \E[\rvz^\top \mathbf{u}(\rvz)] .
\end{align}


\subparagraph{SURE 所需的恒等式。}
设 $\tilde{\rvx}=\rvx+\sigma\rvz$，其中 $\rvz \sim \mathcal{N}(\mathbf{0},\rmI_D)$，$\mathbf{a}$ 为具有适当正则性的任意向量函数，
\begin{align}\label{eq:stein-for-sure}
    \E\big[(\tilde{\rvx}-\rvx)^\top \mathbf{a}(\tilde{\rvx})  \big|  \rvx\big]
 =  \sigma^2 \E\big[ \nabla_{\tilde{\rvx}} \cdot\mathbf{a}(\tilde{\rvx})  \big|  \rvx\big]. 
\end{align}
这是通过应用 \Cref{eq:stein-multi} 并利用链式法则得到的。

\paragraph{从条件 MSE 推导 SURE。}
设 $\rmD:\R^D \to \R^D$ 为去噪器，并定义
\[
R(\rmD;\rvx):=\E \left[\|\rmD(\tilde{\rvx})-\rvx\|_2^2| \rvx\right].
\]
围绕 $\tilde{\rvx}$ 展开：
\[
\begin{aligned}
&~R(\rmD;\rvx) \\
=&~\E \left[\|\rmD(\tilde{\rvx})-\tilde{\rvx}\|^2  \big|  \rvx\right]
+2 \E \left[(\rmD(\tilde{\rvx})-\tilde{\rvx})^\top(\tilde{\rvx}-\rvx)  \big|  \rvx\right]
+\E \left[\|\tilde{\rvx}-\rvx\|^2  \big|  \rvx\right] \\
=&~\E \left[\|\rmD(\tilde{\rvx})-\tilde{\rvx}\|^2  \big|  \rvx\right]
+2\Big(\underbrace{\E[(\tilde{\rvx}-\rvx)^\top \rmD(\tilde{\rvx})|\rvx]}_{\substack{\sigma^2 \E[ \nabla_{\tilde{\rvx}} \cdot\rmD(\tilde{\rvx})|\rvx] \\  \text{ by \Cref{eq:stein-for-sure}}
}}
-\underbrace{\E[(\tilde{\rvx}-\rvx)^\top \tilde{\rvx}|\rvx]}_{ \sigma^2 D}\Big)
\\&\qquad\qquad\qquad+\underbrace{\E[\|\tilde{\rvx}-\rvx\|^2|\rvx]}_{ \sigma^2 D} \\
=&~\E \left[\|\rmD(\tilde{\rvx})-\tilde{\rvx}\|^2
+2\sigma^2  \nabla_{\tilde{\rvx}} \cdot\rmD(\tilde{\rvx})
-D\sigma^2 | \rvx\right].
\end{aligned}
\]
因此，\emph{可观测的}代理
\[ 
\mathrm{SURE}(\rmD;\tilde{\rvx})
:= \|\rmD(\tilde{\rvx})-\tilde{\rvx}\|_2^2
+2\sigma^2  \nabla_{\tilde{\rvx}} \cdot\rmD(\tilde{\rvx})
-D\sigma^2 
\]
满足 $\E \left[\mathrm{SURE}(\rmD;\tilde{\rvx}) \big| \rvx\right]=R(\rmD;\rvx)$。
因此，最小化 SURE（期望意义上或经验意义上）等价于在仅使用含噪观测的情况下最小化真实的条件 MSE。



\subsection{定理~\ref{thm:fpe}：通过 Fokker–Planck 实现边缘对齐}\label{app-sec:scoresde-fpe-proof}

\paragraph{\textit{证明。}}
\paragraph{第一部分：前向 SDE 的 Fokker-Planck 方程。}

考虑前向 SDE：
\[
\diff \rvx(t) = \rvf(\rvx(t), t)  \diff t + g(t)  \diff \rvw(t).
\]
扩散矩阵为 $\sigma(t) = g(t) I_D$，故 $\sigma(t) \sigma(t)^T = g^2(t) I_D$。$\rvx(t)$ 的边缘密度 $p_t(\rvx)$ 满足的 Fokker-Planck 方程为：
\[
\partial_t p_t(\rvx) = -\nabla_{\rvx} \cdot \bigl[ \rvf(\rvx, t) p_t(\rvx) \bigr] + \frac{1}{2} \sum_{i,j=1}^D \frac{\partial^2}{\partial x_i \partial x_j} \bigl[ (g^2(t) \delta_{ij}) p_t(\rvx) \bigr].
\]
计算扩散项：
\[
\sum_{i,j=1}^D \frac{\partial^2}{\partial x_i \partial x_j} \bigl[ g^2(t) \delta_{ij} p_t(\rvx) \bigr] = \sum_{i=1}^D \frac{\partial^2}{\partial x_i^2} \bigl[ g^2(t) p_t(\rvx) \bigr] = g^2(t) \Delta_{\rvx} p_t(\rvx).
\]
因此：
\[
\partial_t p_t(\rvx) = -\nabla_{\rvx} \cdot \bigl[ \rvf(\rvx, t) p_t(\rvx) \bigr] + \frac{1}{2} g^2(t) \Delta_{\rvx} p_t(\rvx).
\]
现在，利用以下式子改写：
\[
\tilde \rvf(\rvx, t) = \rvf(\rvx, t) - \frac{1}{2} g^2(t) \nabla_{\rvx} \log p_t(\rvx).
\]
由于 $\nabla_{\rvx} \log p_t(\rvx) = \frac{\nabla_{\rvx} p_t(\rvx)}{p_t(\rvx)}$，计算：
\[
\nabla_{\rvx} \cdot \bigl[ \tilde \rvf(\rvx, t) p_t(\rvx) \bigr] = \nabla_{\rvx} \cdot \biggl[ \rvf(\rvx, t) p_t(\rvx) - \frac{1}{2} g^2(t) \frac{\nabla_{\rvx} p_t(\rvx)}{p_t(\rvx)} p_t(\rvx) \biggr].
\]
第二项为：
\[
\nabla_{\rvx} \cdot \biggl[ -\frac{1}{2} g^2(t) \nabla_{\rvx} p_t(\rvx) \biggr] = -\frac{1}{2} g^2(t) \Delta_{\rvx} p_t(\rvx).
\]
因此：
\[
\nabla_{\rvx} \cdot \bigl[ \tilde \rvf(\rvx, t) p_t(\rvx) \bigr] = \nabla_{\rvx} \cdot \bigl[ \rvf(\rvx, t) p_t(\rvx) \bigr] - \frac{1}{2} g^2(t) \Delta_{\rvx} p_t(\rvx).
\]
因此：
\[
\partial_t p_t(\rvx) = -\nabla_{\rvx} \cdot \bigl[ \tilde \rvf(\rvx, t) p_t(\rvx) \bigr],
\]
这验证了两种形式的 Fokker-Planck 方程。

\paragraph{第二部分：PF-ODE 的边缘密度。}

考虑 PF-ODE：
\[
\frac{\diff \tilde\rvx(t)}{\diff t} = \tilde \rvf(\tilde\rvx(t), t), \quad \tilde \rvf(\rvx, t) = \rvf(\rvx, t) - \frac{1}{2} g^2(t) \nabla_{\rvx} \log p_t(\rvx).
\]

\textbfs{前向：$\tilde\rvx(0) \sim p_0$。 } 设 $\tilde\rvx(t)$ 满足以 $\tilde\rvx(0) \sim p_0$ 为初值的 PF-ODE。流映射 $\bm{\Psi}_t: \mathbb{R}^D \to \mathbb{R}^D$ 定义为：
\[
\frac{\diff}{\diff t} \bm{\Psi}_t(\mathbf{x}_0) = \tilde \rvf(\bm{\Psi}_t(\mathbf{x}_0), t), \quad \bm{\Psi}_0(\mathbf{x}_0) = \mathbf{x}_0.
\]
由于 $\tilde\rvx(t) = \bm{\Psi}_t(\tilde\rvx(0))$，$\tilde\rvx(t)$ 的密度 $\tilde p_t(\rvx)$ 满足连续性方程：
\[
\partial_t \tilde p_t(\rvx) = -\nabla_{\rvx} \cdot \bigl[ \tilde \rvf(\rvx, t) \tilde p_t(\rvx) \bigr].
\]
由于 $\tilde\rvx(0) \sim p_0$，有 $\tilde p_0(\rvx) = p_0(\rvx)$。由第一部分，$p_t(\rvx)$ 满足：
\[
\partial_t p_t(\rvx) = -\nabla_{\rvx} \cdot \bigl[ \tilde \rvf(\rvx, t) p_t(\rvx) \bigr].
\]
$\tilde p_t$ 和 $p_t$ 满足相同的连续性方程且有相同的初始条件 $p_0$。在足够光滑（例如 $\tilde \rvf \in C^1$）的假设下，解在某个适当的函数空间中是唯一的，故 $\tilde p_t = p_t$。因此对所有 $t \in [0, T]$，$\tilde\rvx(t) \sim p_t$。

\vspace{0.3cm}
\textbfs{后向：$\tilde\rvx(T) \sim p_T$。 } 现在，设 $\tilde\rvx(t)$ 从 $t = T$ 到 $t = 0$ 反向满足 PF-ODE，且 $\tilde\rvx(T) \sim p_T$。该 ODE 为：
\[
\frac{\diff}{\diff t} \tilde\rvx(t) = \tilde \rvf(\tilde\rvx(t), t).
\]
令 $s = T - t$，则反向 ODE 变为：
\[
\frac{\diff}{\diff s} \tilde\rvx(T - s) = -\tilde \rvf(\tilde\rvx(T - s), T - s).
\]
$\tilde\rvx(T - s)$ 的密度 $\tilde p_{T-s}(\rvx)$ 满足：
\[
\partial_s \tilde p_{T-s}(\rvx) = \nabla_{\rvx} \cdot \bigl[ \tilde \rvf(\rvx, T - s) \tilde p_{T-s}(\rvx) \bigr].
\]
由于 $\tilde\rvx(T) \sim p_T$，有 $\tilde p_T = p_T$。$p_t$ 在 $t = T - s$ 处的 Fokker-Planck 方程为：
\[
\partial_t p_{T-s}(\rvx) = -\nabla_{\rvx} \cdot \bigl[ \tilde \rvf(\rvx, T - s) p_{T-s}(\rvx) \bigr].
\]
由于 $\partial_t = -\partial_s$，得：
\[
\partial_s p_{T-s}(\rvx) = \nabla_{\rvx} \cdot \bigl[ \tilde \rvf(\rvx, T - s) p_{T-s}(\rvx) \bigr].
\]
$\tilde p_{T-s}$ 和 $p_{T-s}$ 满足相同的 PDE 且在 $s = 0$ 处有相同的初始条件（$\tilde p_T = p_T$）。由唯一性可得 $\tilde p_{T-s} = p_{T-s}$，故对所有 $t \in [0, T]$，$\tilde\rvx(t) = \tilde\rvx(T - s) \sim p_{T-s} = p_t$。

\paragraph{第三部分：反向时间 SDE 的边缘密度。}
考虑反向时间 SDE：
\[
\diff \bar\rvx(t) = \bigl[ \rvf(\bar\rvx(t), t) - g^2(t) \nabla_{\rvx} \log p_t(\bar\rvx(t)) \bigr] \diff t + g(t) \diff \bar\rvw(t),
\]
其中 $\bar\rvx(0) \sim p_T$，$\bar\rvw(t)$ 是反向时间的标准 Wiener 过程，定义为 $\bar\rvw(t) = \rvw(T - t) - \rvw(T)$。我们需要证明 $\bar\rvx(t) \sim p_{T-t}$。

改写漂移项：
\[
\rvf(\rvx, t) = \tilde \rvf(\rvx, t) + \tfrac{1}{2} g^2(t) \nabla_{\rvx} \log p_t(\rvx),
\]
故：
\[
\rvf(\rvx, t) - g^2(t) \nabla_{\rvx} \log p_t(\rvx)
= \tilde \rvf(\rvx, t) - \tfrac{1}{2} g^2(t) \nabla_{\rvx} \log p_t(\rvx).
\]
因此反向时间 SDE 为
\[
\diff \bar\rvx(t)
=
\biggl[
\tilde \rvf(\bar\rvx(t), t)
- \tfrac{1}{2} g^2(t)\nabla_{\rvx}\log p_t(\bar\rvx(t))
\biggr]\diff t
+ g(t)\diff \bar\rvw(t).
\]

令 $s = T - t$，则 $\diff t = -\diff s$。于是
\begin{align*}
\diff \bar\rvx(T - s)
&=
\biggl[
-\tilde \rvf(\bar\rvx(T - s), T - s)
+ \tfrac{1}{2} g^2(T - s)\nabla_{\rvx}\log p_{T-s}(\bar\rvx(T - s))
\biggr]\diff s \\
&\qquad\qquad\qquad\qquad
+ g(T - s)\diff \bar\rvw(T - s).
\end{align*}
由于 $\bar\rvw(t) = \rvw(T - t) - \rvw(T)$，有
\[
\diff \bar\rvw(T - s) = -\diff \rvw(s).
\]
因此，定义 $\bar\rvw'(s) := -\rvw(s)$（它仍然是标准 Wiener 过程），并将 $\bar\rvx(T-s)$ 重新标记为 $\bar\rvx(s)$，我们得到
\[
\diff \bar\rvx(s)
=
\biggl[
-\tilde \rvf(\bar\rvx(s), T - s)
+ \tfrac{1}{2} g^2(T - s)\nabla_{\rvx}\log p_{T-s}(\bar\rvx(s))
\biggr]\diff s
+ g(T - s)\diff \bar\rvw'(s).
\]

记 $\bar p_s(\rvx)$ 为 $\bar\rvx(s)$ 的密度。其 Fokker-Planck 方程为
\begin{align*}
\partial_s \bar p_s(\rvx)
&=
-\nabla_{\rvx}\cdot
\biggl[
\biggl(
-\tilde \rvf(\rvx, T - s)
+ \tfrac{1}{2} g^2(T - s)\nabla_{\rvx}\log p_{T-s}(\rvx)
\biggr)\bar p_s(\rvx)
\biggr] \\
&\qquad\qquad\qquad\qquad
+ \tfrac{1}{2} g^2(T - s)\Delta_{\rvx}\bar p_s(\rvx).
\end{align*}

假设 $\bar p_s = p_{T-s}$。由第一部分，$p_{T-s}$ 的 Fokker-Planck 方程为
\[
\partial_t p_{T-s}(\rvx)
=
-\nabla_{\rvx}\cdot\bigl[\tilde \rvf(\rvx, T - s)p_{T-s}(\rvx)\bigr].
\]
由于 $\partial_t = -\partial_s$，得
\[
\partial_s p_{T-s}(\rvx)
=
\nabla_{\rvx}\cdot\bigl[\tilde \rvf(\rvx, T - s)p_{T-s}(\rvx)\bigr].
\]

现在将 $\bar p_s = p_{T-s}$ 代入 Fokker-Planck 方程：
\begin{align*}
\partial_s p_{T-s}(\rvx)
&=
-\nabla_{\rvx}\cdot
\biggl[
-\tilde \rvf(\rvx, T - s)p_{T-s}(\rvx)
+ \tfrac{1}{2} g^2(T - s)\frac{\nabla_{\rvx}p_{T-s}(\rvx)}{p_{T-s}(\rvx)}\,p_{T-s}(\rvx)
\biggr] \\
&\qquad\qquad\qquad\qquad
+ \tfrac{1}{2} g^2(T - s)\Delta_{\rvx}p_{T-s}(\rvx) \\
&=
\nabla_{\rvx}\cdot\bigl[\tilde \rvf(\rvx, T - s)p_{T-s}(\rvx)\bigr]
-\tfrac{1}{2} g^2(T - s)\Delta_{\rvx}p_{T-s}(\rvx)
+\tfrac{1}{2} g^2(T - s)\Delta_{\rvx}p_{T-s}(\rvx) \\
&=
\nabla_{\rvx}\cdot\bigl[\tilde \rvf(\rvx, T - s)p_{T-s}(\rvx)\bigr].
\end{align*}
因此，$\bar p_s = p_{T-s}$ 满足 Fokker-Planck 方程。由于 $\bar\rvx(0) \sim p_T$，有 $\bar p_0 = p_T$，与初始条件吻合。唯一性（在足够光滑的条件下）保证了 $\bar p_s = p_{T-s}$，故 $\bar\rvx(t) \sim p_{T-t}$。
\hfill$\blacksquare$


\subsection{命题~\ref{dsm-minimizer}：SM 与 DSM 的最小化子}\label{app-sec:min-dsm}

\paragraph{\textit{证明。}} 为求最小化子 $\mathbf{s}^*$，我们首先考虑固定时间 $t$ 并分析目标函数中的内层期望：
\begin{align*}
\mathcal  J(t, \bm{\phi}) 
&:= \mathbb{E}_{\rvx_0 \sim p_{\text{data}}} 
\mathbb{E}_{\rvx_t \sim p_t(\cdot|\rvx_0)} 
\left[\left\| \mathbf{s}_{\bm{\phi}}(\rvx_t, t) - \nabla_{\rvx_t} \log p_t(\rvx_t|\rvx_0) \right\|_2^2 \right].
\end{align*}
为使该期望最小化，我们需要找到
$\mathbf{s}_{\bm{\phi}}(\rvx_t, t)$
使每个 $\rvx_t$ 的期望平方误差最小化。
我们可以利用 $X_0$ 和 $X_t$ 的联合分布改写该期望：
\begin{align*}
\mathcal J(t, \bm{\phi}) 
= \iint p_{\text{data}}(\rvx_0) p_t(\rvx_t |\rvx_0) 
\left\| \mathbf{s}_{\bm{\phi}}(\rvx_t, t) - \nabla_{\rvx_t} \log p(\rvx_t |  \rvx_0) \right\|_2^2 
 \diff\rvx_0  \diff\rvx_t.
\end{align*}
对每个固定的 $\rvx_t$，我们需要最小化：
\begin{align*}
\int p(\rvx_0 | X_t = \rvx_t) p_t( \rvx_t) 
\left\| \mathbf{s}_{\bm{\phi}}(\rvx_t, t) - \nabla_{\rvx_t} \log p(\rvx_t |  \rvx_0) \right\|_2^2 
 \diff\rvx_0.
\end{align*}
由于 $p_t(\rvx_t)$ 关于
$\mathbf{s}_{\bm{\phi}}(\rvx_t, t)$ 是常数，这等价于最小化：
\begin{align*}
\int p( \rvx_0 | X_t = \rvx_t) 
\left\| \mathbf{s}_{\bm{\phi}}(\rvx_t, t) - \nabla_{\rvx_t} \log p(\rvx_t |  \rvx_0) \right\|_2^2 
  \diff\rvx_0
\end{align*}
当 $\mathbf{s}_{\bm{\phi}}(\rvx_t, t)$ 等于条件期望时，该项最小化：
\begin{align*}
\mathbf{s}^*(\rvx_t, t) 
= \mathbb{E}_{X_0 \sim p(X_0 | X_t = \rvx_t)} 
\left[ \nabla_{\rvx_t} \log p(\rvx_t | X_0) \right].
\end{align*}
现在我们需要证明它等于
$\nabla_{\rvx_t} \log p_t(\rvx_t)$。根据贝叶斯法则和边缘概率的定义：
\begin{align*}
p_t(\rvx_t) 
= \int p_t(\rvx_t | \rvx_0) 
p_{\text{data}}(\rvx_0)   \diff\rvx_0.
\end{align*}
先取对数，再对 $\rvx_t$ 求梯度：
\begin{align*}
\nabla_{\rvx_t} \log p_t(\rvx_t) 
= \frac{\nabla_{\rvx_t} p_t(\rvx_t)}{p_t(\rvx_t)}
= \frac{\nabla_{\rvx_t} \int p_t(\rvx_t | \rvx_0) 
p_{\text{data}}(\rvx_0)   \diff\rvx_0}
{\int p_t(\rvx_t | \rvx_0) 
p_{\text{data}}(\rvx_0)   \diff\rvx_0}.
\end{align*}
在适当的正则性条件下，我们可以交换梯度与积分次序：
\begin{align*}
\nabla_{\rvx_t} \log p_t(\rvx_t) 
= \frac{\int \nabla_{\rvx_t} p_t(\rvx_t | \rvx_0) 
p_{\text{data}}(\rvx_0)   \diff\rvx_0}
{\int p_t(\rvx_t | \rvx_0) 
p_{\text{data}}(\rvx_0)   \diff\rvx_0}.
\end{align*}
\hfill$\blacksquare$



\subsection{高斯的闭式分数函数}

为便于后续引用，我们将一般多元正态分布的分数公式总结为如下引理：

\lemp{高斯的分数}{score-general-gaussian}{设 $\rvx \in \mathbb{R}^D$，考虑多元正态分布
\[
p(\tilde{\rvx} | \rvx) := \mathcal{N}(\tilde{\rvx}; \bm{\mu}, \bm{\Sigma}),
\]
其中 $\bm{\mu} \in \mathbb{R}^D$ 为均值，$\bm{\Sigma} \in \mathbb{R}^{D \times D}$ 为可逆协方差矩阵。其分数函数为
\begin{align}\label{eq:score-gaussian-formula}
    \nabla_{\tilde{\rvx}} \log p(\tilde{\rvx} | \rvx) = -\bm{\Sigma}^{-1} (\tilde{\rvx} - \bm{\mu}).
\end{align}
}{ $p(\tilde{\rvx} | \rvx)$ 的密度函数为：
\[
p(\tilde{\rvx} | \rvx) = \frac{1}{(2\pi)^{D/2} |\bm{\Sigma}|^{1/2}} \exp\left( -\frac{1}{2} (\tilde{\rvx} - \bm{\mu})^\top \bm{\Sigma}^{-1} (\tilde{\rvx} - \bm{\mu}) \right).
\]
为计算分数函数，我们首先对 $ p(\tilde{\rvx} | \rvx) $ 取对数：
\[
\log p(\tilde{\rvx} | \rvx) = -\frac{D}{2} \log(2\pi) - \frac{1}{2} \log |\bm{\Sigma}| - \frac{1}{2} (\tilde{\rvx} - \bm{\mu})^\top \bm{\Sigma}^{-1} (\tilde{\rvx} - \bm{\mu}).
\]
现在，我们计算 $ \log p(\tilde{\rvx} | \rvx) $ 关于 $ \tilde{\rvx} $ 的梯度：
\[
\nabla_{\tilde{\rvx}} \log p(\tilde{\rvx} | \rvx) = - \frac{1}{2} \nabla_{\tilde{\rvx}} \left( (\tilde{\rvx} - \bm{\mu})^\top \bm{\Sigma}^{-1} (\tilde{\rvx} - \bm{\mu}) \right).
\]
利用链式法则，得：
\[
\nabla_{\tilde{\rvx}} \left( (\tilde{\rvx} - \bm{\mu})^\top \bm{\Sigma}^{-1} (\tilde{\rvx} - \bm{\mu}) \right) = 2 \bm{\Sigma}^{-1} (\tilde{\rvx} - \bm{\mu}).
\]
因此，分数函数为：
\begin{align}\label{eq:gaussian-score}
    \nabla_{\tilde{\rvx}} \log p(\tilde{\rvx} | \rvx) = - \bm{\Sigma}^{-1} (\tilde{\rvx} - \bm{\mu}).
\end{align}
}






\clearpage
\newpage

\section{流视角}\label{app-sec:flow-proof}

\subsection{引理~\ref{instant-change-of-var}：瞬时变量替换}
\paragraph{\textit{证明。}}

\textbfs{方法一：变量替换公式。 } 我们将 $p(\rvx(t), t)$ 记为 $p_t(\rvx_t)$。从 ODE 离散化
\[
\rvz_{t+\Delta t} = \rvz_t + \Delta t  \rmF(\rvz_t, t),
\]
出发，正规化流的变量替换公式（\Cref{eq:nf-change-of-var-log}）给出
\begin{align*}
    \log p_{t+\Delta t}(\rvz_{t+\Delta t}) 
    &= \log p_t(\rvz_t) - \log \Big| \det\big(\rmI + \Delta t \nabla_{\rvz} \rmF(\rvz_t, t)\big) \Big| \\
    &= \log p_t(\rvz_t) - \Tr\Big(\log(\rmI + \Delta t \nabla_{\rvz} \rmF(\rvz_t, t)) \Big) \\
    &= \log p_t(\rvz_t) - \Delta t  \Tr\big(\nabla_{\rvz} \rmF(\rvz_t, t)\big) + \mathcal{O}(\Delta t^2),
\end{align*}
其中我们用到了 $\log \det = \Tr \log$ 以及 $\Delta t$ 很小时的展开。
取极限 $\Delta t \to 0$ 即得到对数密度的连续时间微分方程。事实上，\Cref{eq:log-det-tr} 中也应用了同样的技巧。

\vspace{0.3cm}
\textbfs{方法二：连续性方程。 } 我们也可以利用连续性方程，它本质上充当变量替换公式：
\[
\partial_t p(\rvz,t) = - \nabla_{\rvz} \cdot \big(\rmF(\rvz,t) p(\rvz,t)\big).
\]
展开散度，
\[
\partial_t p = - \big( (\nabla_{\rvz} \cdot \rmF) p + \rmF \cdot \nabla_{\rvz} p \big).
\]
沿满足 $\frac{\diff \rvz}{\diff t} = \rmF(\rvz(t), t)$ 的轨迹 $\rvz(t)$，全时间导数为
\begin{align*}
\frac{\diff}{\diff t} p(\rvz(t), t) 
&= \nabla_{\rvz} p \cdot \frac{\diff \rvz}{\diff t} + \partial_t p \\
&= \nabla_{\rvz} p \cdot \rmF - \big( (\nabla_{\rvz} \cdot \rmF) p + \rmF \cdot \nabla_{\rvz} p \big) \\
&= - (\nabla_{\rvz} \cdot \rmF) p.
\end{align*}
两边除以 $p(\rvz(t), t)$，我们得到
\[
\frac{\diff}{\diff t} \log p(\rvz(t), t) = - \nabla_{\rvz} \cdot \rmF(\rvz(t), t).
\]
\hfill$\blacksquare$


\subsection{定理~\ref{thm:continuity-mass}：质量守恒准则}\label{app-sec:mass-conservation}

\paragraph{一些预备知识：流映射与流诱导密度。} 对任意初始位置 $\mathbf{x}_0\in\mathbb{R}^D$，流映射
$\bm{\Psi}_t:\mathbb{R}^D\to\mathbb{R}^D$ 是如下 ODE 的唯一解
\[
\frac{\diff}{\diff t} \bm{\Psi}_t(\mathbf{x}_0)
= \mathbf{v}_t\bigl(\bm{\Psi}_t(\mathbf{x}_0)\bigr), 
\quad \bm{\Psi}_0(\mathbf{x}_0)=\mathbf{x}_0.
\]
在我们的正则性假设下，$\bm{\Psi}_t$ 关于 $t$ 和 $\mathbf{x}_0$ 都是连续可微的。

流诱导密度 $p_t^{\mathrm{fwd}}$ 是初始密度 $p_0$ 在 $\bm{\Psi}_t$ 下的前推：
\[
p_t^{\mathrm{fwd}}(\mathbf{x})
= p_0\bigl(\bm{\Psi}_t^{-1}(\mathbf{x})\bigr)
  \Bigl|\det \bigl(\nabla\bm{\Psi}_t^{-1}(\mathbf{x})\bigr)\Bigr|.
\]
这给出了从 $p_0=p_{\mathrm{data}}$ 出发、并在速度场 $\mathbf{v}_t$ 下演化的粒子在位置 $\mathbf{x}$、时刻 $t$ 处的密度。


\paragraph{\textit{非严格证明：}} \textbfs{充分条件：$p_t^{\mathrm{fwd}} = p_t$ $\Rightarrow$ 连续性方程。 }
在 \Cref{subsec:cov-continuity-eq} 中，我们通过对离散变量替换公式取连续时间极限，得到了连续性方程的\emph{强解}。在那种方法中，假设密度 $p_t$ 足够光滑，使得所有导数经典地存在且 PDE 逐点成立。这里，我们给出一种\emph{弱意义}下的互补推导：连续性方程仅在积分到任意光滑测试函数之后才施加，这放宽了对 $p_t$ 和速度场 $\mathbf{v}_t$ 的正则性要求。这种弱形式不仅因能容纳正则性更低的解而更为严格，而且是 PDE 理论和数值分析中的标准框架。


对任意具有紧支撑的光滑测试函数 $ \varphi(\mathbf{x}) $，前推性质给出：
\begin{align*}
    \int p_t^{\mathrm{fwd}}(\mathbf{x}) \varphi(\mathbf{x})  \diff\mathbf{x} &= \int p_0(\bm{\Psi}_t^{-1}(\mathbf{x})) \left| \det\left( \nabla \bm{\Psi}_t^{-1}(\mathbf{x}) \right) \right| \varphi(\mathbf{x}) \diff\mathbf{x}\\ &=  \int p_0(\mathbf{y}) \varphi(\bm{\Psi}_t(\mathbf{y})) \diff\mathbf{y},
\end{align*}
其中第二个等式由变量替换 $ \mathbf{x} = \bm{\Psi}_t(\mathbf{y}) $ 得到，且 $ \diff\mathbf{y} = \left| \det\left( \nabla \bm{\Psi}_t^{-1}(\mathbf{x}) \right) \right| \diff\mathbf{x} $。


两边对 $ t $ 求导：
\[
\frac{\diff}{\diff t} \int p_t^{\mathrm{fwd}}(\mathbf{x}) \varphi(\mathbf{x}) \diff\mathbf{x} = \frac{\diff}{\diff t} \int p_0(\mathbf{y}) \varphi(\bm{\Psi}_t(\mathbf{y})) \diff\mathbf{y}.
\]

左端为：
\[
\int \frac{\partial p_t^{\mathrm{fwd}}}{\partial t}(\mathbf{x}) \varphi(\mathbf{x}) \diff\mathbf{x}.
\]

右端为：
\[
\int p_0(\mathbf{y}) \nabla \varphi(\bm{\Psi}_t(\mathbf{y})) \cdot \mathbf{v}_t(\bm{\Psi}_t(\mathbf{y})) \diff\mathbf{y},
\]
因为 $ \frac{\partial \bm{\Psi}_t}{\partial t}(\mathbf{y}) = \mathbf{v}_t(\bm{\Psi}_t(\mathbf{y})) $。
作变量替换 $ \mathbf{x} = \bm{\Psi}_t(\mathbf{y}) $，故
\[
\diff\mathbf{y} = \left| \det\left( \nabla \bm{\Psi}_t^{-1}(\mathbf{x}) \right) \right| \diff\mathbf{x} ,
\]
且
\[
p_0(\mathbf{y}) = p_t^{\mathrm{fwd}}(\mathbf{x}) \left| \det\left( \nabla \bm{\Psi}_t(\mathbf{y}) \right) \right| = \frac{p_t^{\mathrm{fwd}}(\mathbf{x})}{\left| \det\left( \nabla \bm{\Psi}_t^{-1}(\mathbf{x}) \right) \right|}.
\]
因此，右端变为：
\[
\int p_t^{\mathrm{fwd}}(\mathbf{x}) \nabla \varphi(\mathbf{x}) \cdot \mathbf{v}_t(\mathbf{x}) \diff\mathbf{x}.
\]

利用 $ \varphi $ 的紧支撑进行分部积分：
\[
\int p_t^{\mathrm{fwd}}(\mathbf{x}) \nabla \varphi(\mathbf{x}) \cdot \mathbf{v}_t(\mathbf{x}) \diff\mathbf{x} = -\int \varphi(\mathbf{x}) \nabla \cdot (p_t^{\mathrm{fwd}}(\mathbf{x}) \mathbf{v}_t(\mathbf{x})) \diff\mathbf{x}.
\]
令两边相等：
\[
\int \left[ \frac{\partial p_t^{\mathrm{fwd}}}{\partial t}(\mathbf{x}) + \nabla \cdot (p_t^{\mathrm{fwd}}(\mathbf{x}) \mathbf{v}_t(\mathbf{x})) \right] \varphi(\mathbf{x}) \diff\mathbf{x} = 0.
\]
由于 $ \varphi $ 是任意的，我们得到：
\[
\frac{\partial p_t^{\mathrm{fwd}}}{\partial t} + \nabla \cdot (p_t^{\mathrm{fwd}} \mathbf{v}_t) = 0.
\]
给定 $ p_t^{\mathrm{fwd}} = p_t $，这蕴含：
\[
\frac{\partial p_t}{\partial t} + \nabla \cdot (p_t \mathbf{v}_t) = 0.
\]


\textbfs{必要条件：连续性方程 $ \Rightarrow p_t^{\mathrm{fwd}} = p_t $。} 假设 $ p_t $ 满足连续性方程：
\[
\frac{\partial p_t}{\partial t} + \nabla \cdot (p_t \mathbf{v}_t) = 0,
\]
初始条件为 $ p_0(\mathbf{x}) = p_{\mathrm{data}}(\mathbf{x}) $。如上所示，我们知道 $ p_t^{\mathrm{fwd}} $ 满足同样的连续性方程，且：
\[
p_0^{\mathrm{fwd}}(\mathbf{x}) = p_0(\bm{\Psi}_0^{-1}(\mathbf{x})) \left| \det\left( \nabla \bm{\Psi}_0^{-1}(\mathbf{x}) \right) \right| = p_0(\mathbf{x}),
\]
因为 $ \bm{\Psi}_0(\mathbf{x}) = \mathbf{x} $。因此，两个密度共享相同的初始条件 $ p_0 = p_{\mathrm{data}} $。

连续性方程可以改写为：
\[
\frac{\partial p}{\partial t} + \mathbf{v}_t \cdot \nabla p + p \nabla \cdot \mathbf{v}_t = 0.
\]
这是一阶线性 PDE。假设 $ \mathbf{v}_t $ 连续可微且全局 Lipschitz，$ p_t $ 足够光滑，则特征线法保证在光滑函数空间中存在唯一解。由于 $ p_t $ 和 $ p_t^{\mathrm{fwd}} $ 满足相同的 PDE 和初始条件，我们得到：
\[
p_t(\mathbf{x}) = p_t^{\mathrm{fwd}}(\mathbf{x})
\]
对所有 $t\in[0,1]$ 和 $\mathbf{x} \in \mathbb{R}^D$ 成立。

这就完成了等价性的证明。
\hfill$\blacksquare$

\clearpage
\newpage

\section{理论补充：扩散模型的统一系统视角}\label{app-sec:equiv-para}

\subsection{命题~\ref{equi-para}：参数化的等价性}

\paragraph{\textit{证明：}} 
正如关于 DSM 损失的定理~\ref{dsm-minimizer} 所述，匹配损失的全局最优
\[
\mathbb{E}_t\left[\omega(t)\,\mathbb{E}_{\rvx_0,\bm{\epsilon}}\left[\|\cdot\|_2^2\right]\right]
\]
在如下内层期望
\[
\mathbb{E}_{\rvx_0,\bm{\epsilon}}\left[\|\cdot\|_2^2\right]
\]
对每个固定的 $t$ 取得最小值时达到。由于这是一个标准的均方误差问题，最小化子是唯一的。
根据去噪分数匹配~\citep{vincent2011connection}，定理~\ref{dsm-minimizer} 表明最优分数函数满足
\begin{align*}
    \rvs^*(\rvx_t,t) = \mathbb{E}_{p(\rvx_0|\rvx_t)}\left[\nabla_{\rvx}\log p_t(\rvx_t|\rvx_0)\right] = \nabla_{\rvx_t} \log p_t(\rvx_t).
\end{align*}
对 $\bm{\epsilon} \sim \mathcal{N}(\bm{0},\rmI)$，利用恒等式 $\nabla_{\rvx}\log p_t(\rvx_t|\rvx_0) = -\frac{1}{\sigma_t} \bm{\epsilon}$，得
\begin{align*}
    \rvs^*(\rvx_t,t) = -\frac{1}{\sigma_t}\, \bm{\epsilon}^*(\rvx_t,t),
\end{align*}
其中 $\bm{\epsilon}^*(\rvx_t,t) = \mathbb{E}_{p(\rvx_0|\rvx_t)}[\bm{\epsilon}|\rvx_t]$ 是 $\mathcal{L}_{\text{noise}}({\bm{\phi}})$ 的最优 $\beps$-预测。
对于 $\rvx$-预测损失 $\mathcal{L}_{\text{clean}}$，最优估计子为
\[
\rvx^*(\rvx_t,t) = \mathbb{E}_{p(\rvx_0|\rvx_t)}[\rvx_0|\rvx_t],
\]
根据 Tweedie 公式，它与分数的关系为
\[
\alpha_t\, \rvx^*(\rvx_t,t) = \rvx_t + \sigma_t^2\, \rvs^*(\rvx_t,t).
\]
对于 $\rvv$-预测损失 $\mathcal{L}_{\text{velocity}}$，最优估计子为
\begin{align*}
    \rvv^*(\rvx_t,t) &= \mathbb{E}_{p(\rvx_0|\rvx_t)}[\alpha_t' \rvx_0 + \sigma_t' \bm{\epsilon}|\rvx_t] \\
    &= \alpha_t' \rvx^* + \sigma_t' \bm{\epsilon}^*.
\end{align*}
将 $\rvx^*$ 和 $\bm{\epsilon}^*$ 用 $\rvs^*$ 表示的表达式代入，得
\begin{align*}
    \rvv^*(\rvx_t,t)
    &= \frac{\alpha_t'}{\alpha_t} \rvx_t + \left( \frac{\alpha_t'}{\alpha_t} \sigma_t^2 - \sigma_t \sigma_t' \right) \rvs^*(\rvx_t,t) \\
    &= f(t) \rvx_t - \frac{1}{2}g^2(t) \rvs^*(\rvx_t,t),
\end{align*}
推导完毕。


\hfill$\blacksquare$



\clearpage
\newpage


\section{理论补充：学习快速扩散生成器}\label{app-sec:flow-map}

\subsection{知识蒸馏损失作为通用框架 \Cref{eq:general-cm-loss} 的一个实例}

我们从 oracle KD 损失出发
\[
\mathcal{L}_{\mathrm{KD}}^{\text{oracle}}(\btheta)
=\E_{\rvx_T\sim p_T}\,
\big\|\rmG_\btheta(\rvx_T,T,0)-\bPsi_{T\to 0}(\rvx_T)\big\|_2^2,
\]
其中 $p_T = p_{\mathrm{prior}}$。对于确定性 ODE 流映射 $\bPsi$（它满足半群性质且沿轨迹是双射的），边缘分布满足前推恒等式
\[
p_t
= \bPsi_{0\to t}\,\sharp\,p_{\mathrm{data}}
= \bPsi_{T\to t}\,\sharp\,p_{\mathrm{prior}}.
\]
因此，
\[
\bPsi_{s\to T}\,\sharp\,p_s=p_T
\quad\text{and}\quad
\bPsi_{T\to 0}\circ\bPsi_{s\to T}=\bPsi_{s\to 0}.
\]
作变量替换 $\rvx_T=\bPsi_{s\to T}(\rvx_s)$，其中 $\rvx_s\sim p_s$，得
\[
\mathcal{L}_{\mathrm{KD}}^{\text{oracle}}(\btheta)
= \E_{\rvx_s\sim p_s} 
\Big\|\rmG_\btheta \big(\bPsi_{s\to T}(\rvx_s),T,0\big)
- \bPsi_{s\to 0}(\rvx_s)\Big\|_2^2.
\]
定义后拉学生 $\widetilde{\rmG}_\btheta(\rvx_s,s,0):=
\rmG_\btheta(\bPsi_{s\to T}(\rvx_s),T,0)$，以将同一损失表示为统一的流映射形式（在 $t=0$ 处）：
\begin{align*}
\mathcal{L}_{\mathrm{KD}}^{\text{oracle}}(\btheta)
= \E_{\rvx_s\sim p_s} 
\big\|\widetilde{\rmG}_\btheta(\rvx_s,s,0)-\bPsi_{s\to 0}(\rvx_s)\big\|_2^2.
\end{align*}
该推导依赖于通过 oracle 流的变量替换以及半群性质。

\hfill$\blacksquare$

\subsection{命题~\ref{vsd-grad} 的参数-流解释}
从命题~\ref{vsd-grad} 的推导中，我们可以将 VSD 的梯度解释为参数诱导的传输流，其中调整模型参数会在数据空间中移动粒子，使其运动与学生和教师分布之间的分数失配对齐。

设 $t\sim p(t)$，$\rvz\sim p(\rvz)$，$\beps\sim\mathcal N(\bm{0},\rmI)$，且
\[
\hat\rvx_t=\alpha_t\,\rmG_\btheta(\rvz)+\sigma_t\,\beps.
\]
定义\emph{样本（粒子）速度}
\[
\widehat{\rvv}_\btheta(\hat\rvx_t):=\partial_\btheta \hat\rvx_t=\alpha_t\,\partial_\btheta \rmG_\btheta(\rvz),
\]
并将 $\rvx$ 空间中的\emph{速度场}定义为条件平均
\[
\rvv_\btheta(\rvx):=\E \big[\widehat{\rvv}_\btheta(\hat\rvx_t) \big| \hat\rvx_t=\rvx\big].
\]

在此定义下，对每个固定的 $t$，密度满足参数-流连续性方程
\[
\partial_\btheta p_t^{\btheta}(\rvx)+\nabla_{\rvx}\cdot\big(p_t^{\btheta}(\rvx) \rvv_\btheta(\rvx)\big)=0.
\]

这里 $\rvv_\btheta(\hat\rvx_t)=\partial_\btheta \hat\rvx_t$ 是
数据空间中（固定 $t$ 时）的\emph{参数诱导的粒子速度}。
等价地，对每个固定的 $t$，密度满足关于 $\btheta$ 的连续性方程：
\[
\partial_\btheta p_t^{\btheta}(\rvx) + \nabla_{\rvx} \cdot \big(p_t^{\btheta}(\rvx) \rvv_\btheta(\rvx)\big) = 0.
\]
因此 $\mathcal{L}_{\text{VSD}}$ 的梯度取传输形式，
\[
\nabla_{\btheta}\mathcal{L}_{\text{VSD}}
= \E\big[\omega(t)\langle
\underbrace{\nabla_{\rvx}\log p_t^{\btheta}-\nabla_{\rvx}\log p_t}_{\text{score mismatch at fixed }t},
\underbrace{\rvv_\btheta}_{\text{parameter–flow velocity}}\rangle\big],
\]
其含义是：调整参数-流使得粒子运动与局部分数失配对齐，
从而减小沿轨迹的散度。









% \begin{align*}
% &\nabla_{\btheta} \int p_t^{\btheta}(\rvx_t)
% \log\frac{p_t^{\btheta}(\rvx_t)}{p_t(\rvx_t)} \diff\rvx_t
% \\&= \int \partial_{\btheta}p_t^{\btheta}(\rvx_t)
%       \log\frac{p_t^{\btheta}(\rvx_t)}{p_t(\rvx_t)} \diff\rvx_t
%    + \int p_t^{\btheta}(\rvx_t) 
%       \partial_{\btheta}\log\frac{p_t^{\btheta}(\rvx_t)}{p_t(\rvx_t)} \diff\rvx_t \\ 
% &= \int \partial_{\btheta}p_t^{\btheta}(\rvx_t)
%       \log\frac{p_t^{\btheta}(\rvx_t)}{p_t(\rvx_t)} \diff\rvx_t
%    + \int p_t^{\btheta}(\rvx_t) 
%       \partial_{\btheta}\log p_t^{\btheta}(\rvx_t) \diff\rvx_t \\
% &= \int \partial_{\btheta}p_t^{\btheta}(\rvx_t)
%       \log\frac{p_t^{\btheta}(\rvx_t)}{p_t(\rvx_t)} \diff\rvx_t
%    + \int \partial_{\btheta}p_t^{\btheta}(\rvx_t) \diff\rvx_t \\
% &= \int \partial_{\btheta}p_t^{\btheta}(\rvx_t)
%       \log\frac{p_t^{\btheta}(\rvx_t)}{p_t(\rvx_t)} \diff\rvx_t
% \end{align*}
% Last is because $\int \partial_{\btheta}p_t^{\btheta}= \partial_{\btheta}1=0.$
% Reparametrizing trick: 
% \[
% \hat\rvx_t(\rvz,\beps;\btheta) = \alpha_t\rmG_\btheta(\rvz) + \sigma_t \beps
% \]
% with $\beps\sim\mathcal{N}(\bm{0},\rmI)$ and $\rvz\sim\mathcal{N}(\bm{0},\rmI)$.
% Hence, 
% \begin{align*}
%     p_t^{\btheta}(\rvx) =\mathbb E_{\rvz,\beps}\big[\delta\big(\rvx-\hat\rvx_t(\rvz,\beps;\btheta)\big)\big],
% \end{align*}
% Chain rule implies:
% \begin{align*}
%     &\partial_\btheta p_t^{\btheta}(\rvx) =   \iint  \left[\nabla_{\rvx}\delta \big(\rvx-\hat\rvx_t\big) \partial_{\btheta}\hat\rvx_t\right] \diff\rvz \diff\beps.
% \end{align*}
% Plugging this back into the identity:
% \begin{align*}
% &\nabla_{\btheta} \int p_t^{\btheta}(\rvx_t)
% \log\frac{p_t^{\btheta}(\rvx_t)}{p_t(\rvx_t)} \diff\rvx_t
% \\
% &= \int \partial_{\btheta}p_t^{\btheta}(\rvx_t)
%       \log\frac{p_t^{\btheta}(\rvx_t)}{p_t(\rvx_t)} \diff\rvx_t
% \\ &= - \int \iint \left[\nabla_{\rvx}\delta \big(\rvx-\hat\rvx_t\big) \partial_{\btheta}\hat\rvx_t\right] 
%       \log\frac{p_t^{\btheta}(\rvx_t)}{p_t(\rvx_t)} \diff\rvx_t \diff\rvz \diff\beps
% \\ &= \int \iint \delta \big(\rvx-\hat\rvx_t\big) \partial_{\btheta}\hat\rvx_t 
%       \nabla_{\rvx} \log\frac{p_t^{\btheta}(\rvx_t)}{p_t(\rvx_t)} \diff\rvx_t \diff\rvz \diff\beps
% \\ &= \iint \partial_{\btheta}\hat\rvx_t 
%       \nabla_{\rvx} \log\frac{p_t^{\btheta}(\rvx_t)}{p_t(\rvx_t)}  \diff\rvz \diff\beps
% \end{align*}
% where we used IBP in the second last identity.

% \hfill$\blacksquare$



% \subsection{Theorem~\ref{thm:ct-cm}: CM Equals CT up to $\smallO(\Delta s)$}

% \paragraph{\textit{Proof:}} 

% \textbfs{Step 1: DDIM Update (with Oracle Score) Is the Conditional Mean.}
% \begin{align*}
% \hat\rvx_{s'}
% &:= \frac{\alpha_{s'}}{\alpha_s} \rvx_s 
% + \sigma_s^2 \left(\frac{\alpha_{s'}}{\alpha_s}-\frac{\sigma_{s'}}{\sigma_s}\right)\nabla_{\rvx_s}\log p_s(\rvx_s)
% \\
% &= \frac{\alpha_{s'}}{\alpha_s}\bigl(\rvx_s+\sigma_s^2\nabla_{\rvx_s}\log p_s(\rvx_s)\bigr)
% - \sigma_{s'}\sigma_s \nabla_{\rvx_s}\log p_s(\rvx_s)
% \\
% &= \alpha_{s'} \E[\rvx_0|\rvx_s]
% + \sigma_{s'}\Bigl(\frac{\rvx_s-\alpha_s \E[\rvx_0|\rvx_s]}{\sigma_s}\Bigr)
% \\
% &= \alpha_{s'} \E[\rvx_0|\rvx_s] + \sigma_{s'} \E[\beps|\rvx_s]
% \\
% &= \E[\alpha_{s'}\rvx_0+\sigma_{s'}\beps|\rvx_s]
% \\
% &= \E[\rvx_{s'}|\rvx_s].
% \end{align*}
% Here, we use the Tweedie's formula
% $\E[\rvx_0 |\rvx_s]=\frac{\rvx_s+\sigma_s^2\nabla_{\rvx_s}\log p_s(\rvx_s)}{\alpha_s}$ 
% and 
% $\rvx_s=\alpha_s\rvx_0+\sigma_s\beps$
% in the third and fourth equalities.

% \textbfs{Step 2: Expand CT Around the Conditional Mean $\hat\rvx_{s'}$.}

% \begin{align*}
% \mathcal{L}_{\mathrm{CT}}
% &= \E_{s,\rvx_s}\E_{\rvx_{s'}|\rvx_s} \Big[w(s) d\big(\rvf_{\btheta}(\rvx_s,s), \rvf_{\btheta^-}(\rvx_{s'},s')\big)\Big]
% \\
% &\overset{\text{(1)}}{=} \E_{s,\rvx_s}\E_{\rvx_{s'}|\rvx_s} \Big[
% w(s) d\big(\rvf_{\btheta}(\rvx_s,s), \rvf_{\btheta^-}(\hat\rvx_{s'},s')\big)
% \\[-2pt]
% &\qquad\qquad\qquad\qquad
% + w(s) \partial_2 d\big(\rvf_{\btheta}(\rvx_s,s),\rvf_{\btheta^-}(\hat\rvx_{s'},s')\big)\big[\partial_1\rvf_{\btheta^-}(\hat\rvx_{s'},s')(\rvx_{s'}-\hat\rvx_{s'})\big]
% \\[-2pt]
% &\qquad\qquad\qquad\qquad
% + w(s) \mathcal O \big(\|\rvx_{s'}-\hat\rvx_{s'}\|^2\big)\Big]
% \\
% &\overset{\text{(2)}}{=} \E_{s,\rvx_s} \Big[w(s) d\big(\rvf_{\btheta}(\rvx_s,s), \rvf_{\btheta^-}(\hat\rvx_{s'},s')\big)\Big]
% \\
% &\quad+ \E_{s,\rvx_s} \Big[w(s) \partial_2 d\big(\rvf_{\btheta}(\rvx_s,s),\rvf_{\btheta^-}(\hat\rvx_{s'},s')\big)\big[\partial_1\rvf_{\btheta^-}(\hat\rvx_{s'},s') \E_{\rvx_{s'}|\rvx_s}(\rvx_{s'}-\hat\rvx_{s'})\big]\Big]
% \\
% &\quad+ \E_{s,\rvx_s}\E_{\rvx_{s'}|\rvx_s} \Big[w(s) \mathcal{O} \big(\|\rvx_{s'}-\hat\rvx_{s'}\|^2\big)\Big]
% \\
% &\overset{\text{(3)}}{=} \E_{s,\rvx_s} \Big[w(s) d\big(\rvf_{\btheta}(\rvx_s,s), \rvf_{\btheta^-}(\hat\rvx_{s'},s')\big)\Big]
% + \E_{s,\rvx_s}\E_{\rvx_{s'}|\rvx_s} \Big[w(s) \mathcal{O} \big(\|\rvx_{s'}-\hat\rvx_{s'}\|^2\big)\Big]
% \\
% &= \E_{s,\rvx_s} \Big[w(s) d\big(\rvf_{\btheta}(\rvx_s,s), \rvf_{\btheta^-}(\hat\rvx_{s'},s')\big)\Big]
% + \mathcal{O} \Big(\E_{s,\rvx_s}\E_{\rvx_{s'}|\rvx_s}\|\rvx_{s'}-\hat\rvx_{s'}\|^2\Big)
% \\
% &= \mathcal{L}_{\mathrm{CM}} 
% + \mathcal{O} \big(\E_{s,\rvx_s}\E_{\rvx_{s'}|\rvx_s}\|\rvx_{s'}-\hat\rvx_{s'}\|^2\big)
% \\
% &\overset{\text{(4)}}{=} \mathcal{L}_{\mathrm{CM}} + \smallO (\Delta s)
% \end{align*}

% Here, (1) applies a second-order Taylor expansion of 
% \[
% h(\rvx') := d(\rvf_{\btheta}(\rvx_s,s), \rvf_{\btheta^-}(\rvx',s'))
% \]
% around $\hat\rvx_{s'}$ in its second argument. 
% (2) applies the tower property 
% \[
% \E_{s,\rvx_s}\E_{\rvx_{s'}|\rvx_s}[\cdot]=\E_{s,\rvx_s}[\cdot]
% \]
% and notes that, inside $\E_{\rvx_{s'}|\rvx_s}$, the only $\rvx_{s'}$-dependence is through $(\rvx_{s'}-\hat\rvx_{s'})$, so all other factors are treated as constants and the inner expectations reduce to $\E_{\rvx_{s'}|\rvx_s}(\rvx_{s'}-\hat\rvx_{s'})$ and $\E_{\rvx_{s'}|\rvx_s}\|\rvx_{s'}-\hat\rvx_{s'}\|^2$. (3) uses $\E[\rvx_{s'}-\hat\rvx_{s'}|\rvx_s]=\bm{0}$ 
% because $\hat\rvx_{s'}=\E[\rvx_{s'}|\rvx_s]$.
% (4) uses the linear–Gaussian scheduler, which gives 
% \[
% \E[\|\rvx_{s'}-\hat\rvx_{s'}\|^2|\rvx_s]=\mathcal{O}(\Delta s^2),
% \]
% thus leading to a total remainder of $\smallO(\Delta s)$.


% \hfill$\blacksquare$


\subsection{定理~\ref{thm:ct-cm}：CM 与 CT 相差至多 $\mathcal{O}(\Delta s^2)$}

\paragraph{\textit{证明：}} 

\textbfs{第一步：DDIM 更新（使用 Oracle 分数）是条件均值。}
\begin{align*}
\hat\rvx_{s'}
&:= \frac{\alpha_{s'}}{\alpha_s} \rvx_s 
+ \sigma_s^2 \left(\frac{\alpha_{s'}}{\alpha_s}-\frac{\sigma_{s'}}{\sigma_s}\right)\nabla_{\rvx_s}\log p_s(\rvx_s)
\\
&= \frac{\alpha_{s'}}{\alpha_s}\bigl(\rvx_s+\sigma_s^2\nabla_{\rvx_s}\log p_s(\rvx_s)\bigr)
- \sigma_{s'}\sigma_s \nabla_{\rvx_s}\log p_s(\rvx_s)
\\
&= \alpha_{s'} \E[\rvx_0|\rvx_s]
+ \sigma_{s'}\Bigl(\frac{\rvx_s-\alpha_s \E[\rvx_0|\rvx_s]}{\sigma_s}\Bigr)
\\
&= \alpha_{s'} \E[\rvx_0|\rvx_s] + \sigma_{s'} \E[\beps|\rvx_s]
\\
&= \E[\alpha_{s'}\rvx_0+\sigma_{s'}\beps|\rvx_s]
\\
&= \E[\rvx_{s'}|\rvx_s].
\end{align*}
这里我们使用 Tweedie 公式
\[
\E[\rvx_0 |\rvx_s]=\frac{\rvx_s+\sigma_s^2\nabla_{\rvx_s}\log p_s(\rvx_s)}{\alpha_s},
\]
以及共享噪声耦合
\[
\rvx_s=\alpha_s\rvx_0+\sigma_s\beps,
\qquad
\rvx_{s'}=\alpha_{s'}\rvx_0+\sigma_{s'}\beps
\quad (s'=s-\Delta s),
\]
用于第三和第四个等式。我们还注意到，更精确地说，$\hat\rvx_{s'}=\E[\rvx_{s'}|\rvx_s]$ 应理解为
$\hat\rvx_{s'}=\E[\rvx_{s'}|\rvx_s, s]$，因为一旦采样得到 $s$，$s'=s-\Delta s$ 就被确定性地固定了。




\textbfs{第二步：在条件均值 $\hat\rvx_{s'}$ 附近展开 CT。}

回忆 $s$ 先被采样，然后我们确定性地令 $s':=s-\Delta s$。
\begin{align*}
\mathcal{L}_{\mathrm{CT}}
&= \E_{s,\rvx_s}\E_{\rvx_{s'}|\rvx_s} \Big[w(s) d\big(\rvf_{\btheta}(\rvx_s,s), \rvf_{\btheta^-}(\rvx_{s'},s')\big)\Big]
\\
&\overset{\text{(1)}}{=} \E_{s,\rvx_s}\E_{\rvx_{s'}|\rvx_s} \Big[
w(s) d\big(\rvf_{\btheta}(\rvx_s,s), \rvf_{\btheta^-}(\hat\rvx_{s'},s')\big)
\\[-2pt]
&\qquad\qquad\qquad\qquad
+ w(s) \partial_2 d\big(\rvf_{\btheta}(\rvx_s,s),\rvf_{\btheta^-}(\hat\rvx_{s'},s')\big)\big[\partial_1\rvf_{\btheta^-}(\hat\rvx_{s'},s')(\rvx_{s'}-\hat\rvx_{s'})\big]
\\[-2pt]
&\qquad\qquad\qquad\qquad
+ w(s) \mathcal O \big(\|\rvx_{s'}-\hat\rvx_{s'}\|^2\big)\Big]
\\
&\overset{\text{(2)}}{=} \E_{s,\rvx_s} \Big[w(s) d\big(\rvf_{\btheta}(\rvx_s,s), \rvf_{\btheta^-}(\hat\rvx_{s'},s')\big)\Big]
\\
&\quad+ \E_{s,\rvx_s} \Big[w(s) \partial_2 d\big(\rvf_{\btheta}(\rvx_s,s),\rvf_{\btheta^-}(\hat\rvx_{s'},s')\big)\big[\partial_1\rvf_{\btheta^-}(\hat\rvx_{s'},s') \E_{\rvx_{s'}|\rvx_s}(\rvx_{s'}-\hat\rvx_{s'})\big]\Big]
\\
&\quad+ \E_{s,\rvx_s}\E_{\rvx_{s'}|\rvx_s} \Big[w(s) \mathcal{O} \big(\|\rvx_{s'}-\hat\rvx_{s'}\|^2\big)\Big]
\\
&\overset{\text{(3)}}{=} \E_{s,\rvx_s} \Big[w(s) d\big(\rvf_{\btheta}(\rvx_s,s), \rvf_{\btheta^-}(\hat\rvx_{s'},s')\big)\Big]
+ \E_{s,\rvx_s}\E_{\rvx_{s'}|\rvx_s} \Big[w(s) \mathcal{O} \big(\|\rvx_{s'}-\hat\rvx_{s'}\|^2\big)\Big]
\\
&= \mathcal{L}_{\mathrm{CM}} 
+ \mathcal{O} \big(\E\|\rvx_{s'}-\hat\rvx_{s'}\|^2\big)
\\
&\overset{\text{(4)}}{=} \mathcal{L}_{\mathrm{CM}} + \mathcal{O}(\Delta s^2).
\end{align*}
这里，(1) 对
\[
h(\rvx') := d\big(\rvf_{\btheta}(\rvx_s,s), \rvf_{\btheta^-}(\rvx',s')\big)
\]
在第二个参数上于 $\hat\rvx_{s'}$ 处应用二阶 Taylor 展开。(2) 使用\textit{塔性质}（即全期望律）：
\[
\E_{s,\rvx_s}\E_{\rvx_{s'}|\rvx_s,s}[\cdot]
=\E_{s,\rvx_s}[\cdot],
\]
并注意到在 $\E_{\rvx_{s'}|\rvx_s,s}$ 内部，项
\[
w(s),\ \rvf_{\btheta}(\rvx_s,s),\ \hat\rvx_{s'},\ s'
\]
不依赖于随机变量 $\rvx_{s'}$（因为 $s' = s-\Delta s$ 和 $\hat\rvx_{s'}=\E[\rvx_{s'}|\rvx_s,s]$ 在 $(s,\rvx_s)$ 固定后就固定了）。因此唯一的 $\rvx_{s'}$ 依赖性来自 $(\rvx_{s'}-\hat\rvx_{s'})$，内层期望化为
$\E_{\rvx_{s'}|\rvx_s,s}(\rvx_{s'}-\hat\rvx_{s'})$ 和
$\E_{\rvx_{s'}|\rvx_s,s}\|\rvx_{s'}-\hat\rvx_{s'}\|^2$。(3) 使用
\[
\E[\rvx_{s'}-\hat\rvx_{s'}|\rvx_s,s]=\bm 0
\quad\text{because}\quad
\hat\rvx_{s'}=\E[\rvx_{s'}|\rvx_s,s].
\]
对于 (4)，注意
\[
\E\|\rvx_{s'}-\hat\rvx_{s'}\|^2
=\E\,\Var(\rvx_{s'}|\rvx_s,s)
\le \E\|\rvx_{s'}-\rvx_s\|^2.
\]
此外，
\[
\rvx_{s'}-\rvx_s=(\alpha_{s'}-\alpha_s)\rvx_0+(\sigma_{s'}-\sigma_s)\beps,
\]
且对于光滑的调度器，$\alpha_{s'}-\alpha_s=\mathcal{O}(\Delta s)$，
$\sigma_{s'}-\sigma_s=\mathcal{O}(\Delta s)$，故
\[
\E\|\rvx_{s'}-\rvx_s\|^2=\mathcal{O}(\Delta s^2),
\]
从而
\[
\E\|\rvx_{s'}-\hat\rvx_{s'}\|^2=\mathcal{O}(\Delta s^2).
\]


\hfill$\blacksquare$


\subsection{命题~\ref{continuous-time-ct}：CT 梯度的连续时间极限}

\paragraph{\textit{证明：}}

我们通过将 \Cref{eq:cm-discrete} 写成如下形式来简化记号
\begin{align*}
\mathcal{L}_{\mathrm{CM}}^{\Delta s}(\btheta,\btheta^-)
:=
\E \left[\omega(s)\big\|
\rvf_{\btheta}(\rvx_s,s) -
\rvf_{\btheta^-}\big(\bPsi_{s\to s-\Delta s}(\rvx_s),\, s-\Delta s\big)
\big\|_2^2\right].
\end{align*}
为简化记号，记
\[
\tilde\rvx_{s-\Delta s}:=\bPsi_{s\to s-\Delta s}(\rvx_s).
\]

定义
\[
\delta \rvf
:= \rvf_{\btheta}(\rvx_s,s)-\rvf_{\btheta^-}(\rvx_s,s).
\]
则
\begin{align*}
\big\|\rvf_{\btheta}(\rvx_s, s)- \rvf_{\btheta^-}(\tilde\rvx_{s-\Delta s}, s-\Delta s)\big\|_2^2
=
\Big\|\delta \rvf
+
\Big(\rvf_{\btheta^-}(\rvx_s, s)-\rvf_{\btheta^-}(\tilde\rvx_{s-\Delta s}, s-\Delta s)\Big)\Big\|_2^2.
\end{align*}

在 $(\rvx_s,s)$ 处应用 Taylor 展开得
\begin{align*}
\rvf_{\btheta^-}(\rvx_s,s) -
\rvf_{\btheta^-}(\tilde\rvx_{s-\Delta s}, s-\Delta s)
=
\frac{\diff}{\diff s}\rvf_{\btheta^-}(\rvx_s,s)\,\Delta s
+ \mathcal{O}(|\Delta s|^2),
\end{align*}
其中我们用到了 oracle 流映射的一阶展开，
\[
\rvx_s - \bPsi_{s\to s-\Delta s}(\rvx_s)
= \rvv^*(\rvx_s,s)\,\Delta s + \mathcal{O}(|\Delta s|^2),
\]
以及全微分恒等式
\[
\frac{\diff}{\diff s}\rvf_{\btheta^-}(\rvx_s,s)
=
\partial_s \rvf_{\btheta^-}(\rvx_s,s)
+
\big(\partial_{\rvx}\rvf_{\btheta^-}(\rvx_s,s)\big)\,\rvv^*(\rvx_s,s).
\]

因此，
\begin{align*}
\big\|\rvf_{\btheta}(\rvx_s, s)- \rvf_{\btheta^-}(\tilde\rvx_{s-\Delta s}, s-\Delta s)\big\|_2^2
=
\|\delta \rvf\|_2^2
+
2\,\delta \rvf^\top \frac{\diff}{\diff s}\rvf_{\btheta^-}(\rvx_s,s)\,\Delta s
+
\mathcal{O}(|\Delta s|^2).
\end{align*}

现在对 $\btheta$ 求导。由于 $\btheta^-$ 相对于 $\nabla_{\btheta}$ 被视为固定（即不通过 stop-gradient 目标求梯度），因此不通过 $\rvf_{\btheta^-}$ 求梯度。故
\begin{align*}
\nabla_{\btheta}\mathcal{L}_{\mathrm{CM}}^{\Delta s}(\btheta,\btheta^-)
&=
\nabla_{\btheta}\E\!\left[\omega(s)\|\delta \rvf\|_2^2\right]\\
&\qquad\qquad
+
2\,\E\!\left[
\omega(s)\,
\nabla_{\btheta}\rvf_{\btheta}(\rvx_s,s)^\top
\frac{\diff}{\diff s}\rvf_{\btheta^-}(\rvx_s,s)
\right]\Delta s
+
\mathcal{O}(|\Delta s|^2).
\end{align*}

在目标网络是在线网络的 stop-gradient 副本的比较点处，即 $\btheta=\btheta^-$，有
\[
\delta \rvf
=
\rvf_{\btheta}(\rvx_s,s)-\rvf_{\btheta^-}(\rvx_s,s)
=
0,
\]
因此
\[
\nabla_{\btheta}\E\!\left[\omega(s)\|\delta \rvf\|_2^2\right]=0.
\]
于是，
\begin{align*}
\left.
\nabla_{\btheta}\mathcal{L}_{\mathrm{CM}}^{\Delta s}(\btheta,\btheta^-)
\right|_{\btheta=\btheta^-}
&=
\left.
2\,\E\!\left[
\omega(s)\,
\nabla_{\btheta}\rvf_{\btheta}(\rvx_s,s)^\top
\frac{\diff}{\diff s}\rvf_{\btheta^-}(\rvx_s,s)
\right]
\right|_{\btheta=\btheta^-}\Delta s
+
\mathcal{O}(|\Delta s|^2).
\end{align*}

两边除以 $\Delta s$ 并取极限得
\begin{align*}
\left.
\lim_{\Delta s\to 0}\frac{1}{\Delta s}\nabla_{\btheta}\mathcal{L}_{\mathrm{CM}}^{\Delta s}(\btheta,\btheta^-)
\right|_{\btheta=\btheta^-}
&=
\left.
\nabla_{\btheta}\E\!\left[
2\,\omega(s)\,
\rvf_{\btheta}^\top(\rvx_s,s)\,
\frac{\diff}{\diff s}\rvf_{\btheta^-}(\rvx_s,s)
\right]
\right|_{\btheta=\btheta^-}.
\end{align*}
这就证明了在 $\btheta=\btheta^-$ 处的恒等式。

\hfill$\blacksquare$


\newpage
\section{（可选）阐释扩散模型（EDM）}\label{sec:edm}

我们介绍
$\rvx$-预测模型中神经网络参数化设计的特定准则，如《阐释扩散模型》
（Elucidating Diffusion Models，
EDM）~\citep{karras2022elucidating} 所提出。EDM 提供了一种简单的
配方，使训练过程更易优化且更可靠。
$\rvx$-预测模型表示为依赖时间的跳跃连接与缩放残差的组合
（\Cref{eq:D-parametrization}）。其核心思想是在所有时刻将输入和回归目标归一化到
单位方差，并调整跳跃路径，使得残差误差不随时间演化而被放大。
这一配方
已成为扩散模型实现中广泛采用的默认设置，并自然地
推广到流映射学习，特别是 Consistency Model 系列。




\subsection{神经网络 $\rvx_{\bm{\phi}}$ 设计准则}\label{subsec:D-design-edm}

EDM 考虑 $\rvx$-预测模型的一种参数化（略带滥用记号地记为 $\rvx_{\bm{\phi}}(\rvx, t)$）\footnote{如前所述，四种预测类型都是等价的，均可归结为 $\rvx$-预测情形。EDM 采用这一形式，它既被充分研究过，又自然地与流映射模型生成干净样本的目标相一致。}，形式如下：
\begin{align}\label{eq:D-parametrization}
    \rvx_{\bm{\phi}}(\rvx, t) := c_{\mathrm{skip}}(t) \rvx + c_{\mathrm{out}}(t) \rmF_{\bm{\phi}}\left(c_{\mathrm{in}}(t) \rvx, c_{\mathrm{noise}}(t)\right).
\end{align}
这里，$c_{\mathrm{skip}}(t)$、$c_{\mathrm{out}}(t)$、$c_{\mathrm{in}}(t)$ 和 $c_{\mathrm{noise}}(t)$ 是依赖时间的函数。
它们的选取基于即将介绍的实践考量，旨在提升训练过程中的稳定性和性能。



将其代入 \Cref{eq:general-dm-loss-D}，经过直接的代数运算后目标函数变为\footnote{我们从 $\rvx$-预测扩散损失出发，代入 \Cref{eq:D-parametrization} 给出的参数化：
\begin{align*}
    \begin{aligned}
      &~\mathbb{E}_{\rvx_0, \bm{\epsilon}, t} \left[\omega(t)\norm{ \rvx_{\bm{\phi}}(\rvx_t, t) - \rvx_0}_2^2\right] 
      \\= &~\mathbb{E}_{\rvx_0, \bm{\epsilon}, t} \left[\omega(t)\Big\Vert c_{\mathrm{skip}}(t) \underbrace{(\alpha_t \rvx_0+\sigma_t\beps)}_{\rvx_t} + c_{\mathrm{out}}(t) \rmF_{\bm{\phi}}\left(c_{\mathrm{in}}(t) \rvx_t, c_{\mathrm{noise}}(t)\right) - \rvx_0\Big\Vert_2^2\right]
        \\= &~\mathbb{E}_{\rvx_0, \bm{\epsilon}, t} \left[\omega(t)c_{\mathrm{out}}^2(t)\Big\Vert  \rmF_{\bm{\phi}}\left(c_{\mathrm{in}}(t) \rvx, c_{\mathrm{noise}}(t)\right) - \underbrace{\left(\frac{\left(1-c_{\mathrm{skip}}(t)\alpha_t\right)\rvx_0-c_{\mathrm{skip}}(t)\sigma_t \bm{\epsilon}}{c_{\mathrm{out}}(t)}\right)}_{\rvx_{\mathrm{tgt}}(t)}\Big\Vert_2^2\right]
        \\ = &~\mathrm{\Cref{eq:loss-F}}.
    \end{aligned}
\end{align*}

}：
\begin{align}\label{eq:loss-F}
   \mathbb{E}_{\rvx_0, \bm{\epsilon}, t} \left[\omega(t)c_{\mathrm{out}}^2(t)\norm{ \rmF_{\bm{\phi}}\left(c_{\mathrm{in}}(t)\rvx_t, c_{\mathrm{noise}}(t)\right) - \rvx_{\mathrm{tgt}}(t)}_2^2\right].
\end{align}
这里，回归目标 $\rvx_{\mathrm{tgt}}(t)$ 由下式得到：
\[
\rvx_{\mathrm{tgt}}(t) = \frac{\left(1-c_{\mathrm{skip}}(t)\alpha_t\right)\rvx_0-c_{\mathrm{skip}}(t)\sigma_t \bm{\epsilon}}{c_{\mathrm{out}}(t)}.
\]

通过引入 $p_{\text{data}}$ 的标准差（记为 $\sigma_{\mathrm{d}}$），EDM 提出如下网络参数化设计准则，可称为\emph{单位方差准则}。

\paragraph{输入的单位方差。}
\begin{align*}
    &\text{Var}_{\rvx_0,\bm{\epsilon}}\left[ c_{\mathrm{in}}(t)\rvx_t\right] = 1 \nonumber
    \\ \Longleftrightarrow &~c_{\mathrm{in}}^2(t) \text{Var}_{\rvx_0,\bm{\epsilon}}\left[ \alpha_t \rvx_0 +\sigma_t \bm{\epsilon} \right] = 1 \nonumber
    \\ \Longleftrightarrow &~c_{\mathrm{in}}(t)  = \frac{1}{\sqrt{\sigma_{\mathrm{d}}^2 \alpha_t^2 + \sigma_t^2}},
\end{align*}
取正根。
\paragraph{训练目标的单位方差。}
\begin{align}\label{eq:c-out}
    &\text{Var}_{\rvx_0,\bm{\epsilon}}\left[ \rvx_{\mathrm{tgt}}(t) \right] = 1 \nonumber
    \\ \Longleftrightarrow &~c_{\mathrm{out}}^2(t)  = \left(1-c_{\mathrm{skip}}(t)\alpha_t\right)^2 \sigma_{\mathrm{d}}^2 + c_{\mathrm{skip}}^2(t)\sigma_t^2,
\end{align}
其中数据已中心化（$\E[\rvx_0]=\bm{0}$）。
\paragraph{最小化从 $\rmF_{\bm{\phi}}$ 到 $\rvx_{\bm{\phi}}$ 的误差放大。} EDM 旨在缓解网络学习误差从 $\rmF_{\bm{\phi}}$ 到 $\rvx_{\bm{\phi}}$ 的放大。这通过选择 $c_{\mathrm{skip}}$ 来最小化 $c_{\mathrm{out}}$ 实现：
\[
c_{\mathrm{skip}}^* \in \argmin_{c_{\mathrm{skip}}}  c_{\mathrm{out}}^2.
\]
采用标准方法，通过求解 $\frac{\partial c_{\mathrm{out}}}{\partial c_{\mathrm{skip}}} = 0$ 得到临界点 $c_{\mathrm{skip}}^*$，得
\begin{align*}
    c_{\mathrm{skip}}^*(t) = \frac{\alpha_t \sigma_{\mathrm{d}}^2 }{\alpha_t^2 \sigma_{\mathrm{d}}^2  + \sigma_t^2}. 
\end{align*}
将其代入 \Cref{eq:c-out}，最优值为
\begin{align*}
    c_{\mathrm{out}}^*(t) = \pm \frac{\sigma_t \sigma_{\mathrm{d}} }{\sqrt{\alpha_t^2 \sigma_{\mathrm{d}}^2  + \sigma_t^2}}. 
\end{align*}
按惯例，我们使用非负分支作为输出尺度：
\[
c_{\mathrm{out}}^*(t) = \frac{\sigma_t \sigma_{\text d}}
{\sqrt{\alpha_t^2 \sigma_{\text d}^2 + \sigma_t^2}}
\;\;\;(\ge 0),
\]
它保证 $c_{\mathrm{out}}(0)=0$ 且当 $\sigma_t$ 足够大时 $c_{\mathrm{out}}(t)\to\sigma_{\text d}$，
给出直观的极限 $\rvx_{\bm{\phi}}(\rvx_t,0)\approx \rvx_t$ 和
$\rvx_{\bm{\phi}}(\rvx_t,t)\approx \sigma_{\text d}\,\rmF_{\bm{\phi}}(\cdot)$。
% \dk{If not used in subsequent subsections, do we need to have this section? Don't need to specify the rationale of these selections or even describe what are they. Just $c_{skip}$ is enough. Too much information needed to get to the essense of this chapter.}



我们将这些系数总结如下：
\begin{mdframed}
记 $R_t := \alpha_t^2 \sigma_{\mathrm{d}}^2+ \sigma_t^2$，我们有如下选取：
    \begin{align}\label{eq:edm-coeff}
        \begin{aligned}
            c_{\mathrm{in}}(t)  = \frac{1}{\sqrt{R_t}}, \quad
            c_{\mathrm{skip}}(t) = \frac{\alpha_t \sigma_{\mathrm{d}}^2 }{R_t}, \quad
            c_{\mathrm{out}}(t) =  \frac{\sigma_t \sigma_{\mathrm{d}} }{\sqrt{R_t}}.
        \end{aligned}
    \end{align}
\end{mdframed}

利用 \Cref{eq:edm-coeff}，回归目标 $\rvx_{\mathrm{tgt}}(t)$ 简化为
\[
\rvx_{\mathrm{tgt}}(t) = \frac{1}{\sigma_{\mathrm{d}}}\frac{\sigma_t \rvx_0 - \alpha_t \sigma_{\mathrm{d}}^2\bm{\epsilon} }{\sqrt{R_t}}.
\] 
此外，将这些表达式代入 \Cref{eq:D-parametrization} 和 \Cref{eq:loss-F}，可以同时简化参数化和损失函数，得：
\begin{align*}
    \rvx_{\bm{\phi}}(\rvx, t)=\frac{\alpha_t \sigma_{\mathrm{d}}^2 }{R_t} \rvx + \frac{\sigma_t \sigma_{\mathrm{d}} }{\sqrt{R_t}}\rmF_{\bm{\phi}}\left(\frac{1}{\sqrt{R_t}}\rvx, c_{\mathrm{noise}}(t)\right),
\end{align*}
和
\begin{align*}
   \mathbb{E}_{\rvx_0, \bm{\epsilon}, t} \left[\omega(t)\frac{\sigma_t^2  }{R_t}\norm{ \sigma_{\mathrm{d}} \rmF_{\bm{\phi}}\left(\frac{1}{\sqrt{R_t}}\rvx_t, c_{\mathrm{noise}}(t)\right) - \left(\frac{\sigma_t \rvx_0 - \alpha_t \sigma_{\mathrm{d}}^2\bm{\epsilon} }{\sqrt{R_t}}\right)}_2^2\right].
\end{align*}

在此参数化以及条件 $\alpha_0 \approx 1$ 和 $\sigma_0 \approx 0$ 下，我们观察到
\[
c_{\mathrm{skip}}(0) \approx 1 \quad \text{and} \quad c_{\mathrm{out}}(0) \approx 0.
\] 


\paragraph{$c_{\mathrm{noise}}(t)$ 的选取。}
它为 $\rmF_{\bm{\phi}}$ 提供噪声水平嵌入；噪声水平 $\sigma_t$ 的任意
一一映射（例如 $c_{\mathrm{noise}}(t)=\log \sigma_t$）都是合适的。

\subsection{EDM 的一个常见特例：$\alpha_t = 1$，$\sigma_t = t$}

我们考虑 EDM 中使用的简化噪声调度，其中 $\alpha_t = 1$ 且 $\sigma_t = t$，这也出现在 \Cref{sec:CTM} 中关于 CTM 的讨论里。在此设置下，前向过程变为
\[
\rvx_t = \rvx_0 + t\bm{\epsilon}, \quad \text{with } \rvx_0 \sim p_{\text{data}}, \,\, \bm{\epsilon} \sim \mathcal{N}(\bm{0}, \rmI),
\]
对应的扰动核为
\[
p_t(\rvx_t|\rvx_0) = \mathcal{N}(\rvx_t; \rvx_0, t^2\rmI).
\]
相应地，终端时刻的先验分布设为
\[
p_{\text{prior}}(\rvx_T) := \mathcal{N}(\rvx_T; \bm{0}, T^2 \rmI).
\]
由扰动核诱导的边缘密度由如下卷积给出：
\begin{align*}
    p_t(\rvx) = \int \mathcal{N}(\cdot; \bm{0}, t^2 \rmI)p_{\text{data}}(\rvx_0)\diff \rvx_0.
\end{align*}
% The marginal density induced by the perturbation kernel is given by the convolution:
% \begin{align*}
%     p_t(\rvx) = \int p_t(\rvx|\rvx_0)p_{\text{data}}(\rvx_0)\diff \rvx_0 = \left[p_{\text{data}}(\cdot) * \mathcal{N}(\cdot; \bm{0}, t^2 \rmI)\right](\rvx),
% \end{align*}
% where $*$ denotes convolution over densities\dk{Either provide formal definition of it (in appendix) or cite some reference}.

在此设置下，基于 $\rvx$-预测 $\rvx_{\bm{\phi}^\times}$（见 \Cref{eq:clean_pf_ode}）的 PF-ODE 简化为
\begin{align*}
    \frac{\diff \rvx(t)}{\diff t} = \frac{\rvx(t) - \rvx_{\bm{\phi}^\times}(\rvx(t), t)}{t}.
\end{align*}

将此形式代入 \Cref{eq:edm-coeff}，神经网络参数化系数变为
\begin{align}\label{eq:edm-simple-coefficients}
    c_{\mathrm{in}}(t) = \frac{1}{\sqrt{\sigma_{\mathrm{d}}^2 + t^2}}, \quad
    c_{\mathrm{skip}}(t) = \frac{\sigma_{\mathrm{d}}^2}{\sigma_{\mathrm{d}}^2 + t^2}, \quad
    c_{\mathrm{out}}(t) = \pm \frac{t \sigma_{\mathrm{d}}}{\sqrt{\sigma_{\mathrm{d}}^2 + t^2}}.
\end{align}
由这些表达式，我们观察到：
\begin{itemize}
    \item 当 $t \approx 0$ 时，噪声水平可忽略，故
    $c_{\mathrm{skip}} \approx 1$ 且 $c_{\mathrm{out}} \approx 0$。在此极限下，
    跳跃路径占主导，网络本质上直接通过其
    输入，
    \[
    \rvx_{\bm{\phi}}(\rvx, t) \approx \rvx.
    \]
    \item 当 $t \gg 0$ 时，输入被噪声严重破坏，故
    $c_{\mathrm{skip}} \approx 0$ 且 $c_{\mathrm{out}} \approx \sigma_{\mathrm{d}}$。
    在此情形下，跳跃路径消失，模型输出完全
    由学习到的残差决定，
    \[
    \rvx_{\bm{\phi}}(\rvx, t) \approx \sigma_{\mathrm{d}} \rmF_{\bm{\phi}}\left(c_{\mathrm{in}}(t)\rvx, c_{\mathrm{noise}}(t)\right),
    \]
    这意味着网络 $\rmF_{\bm{\phi}}$ 从归一化的含噪输入中预测干净信号的缩放代理；
    在高噪声水平下，模型输出因此
    完全由学习到的去噪函数决定。
\end{itemize}
简言之，该参数化在小 $t$ 时的恒等映射
与大 $t$ 时在标准化输入上的缩放残差预测子之间平滑过渡。
